% Lesson 24 — Writing: Writes compositions on taking part and managing time in community/workplace/school activities — Anglais Tle A — Cours signé OnBuch+
% Programme officiel MINESEC. Compiler avec : tectonic EN24-writing-writes-compositions-on-taking-part-and-managing.tex
\newcommand{\DOCMATIERE}{Anglais}
\newcommand{\DOCNIVEAU}{Tle A}
\newcommand{\DOCMODULE}{Chapter 2 : Economic Life and Occupations}
\newcommand{\DOCLECON}{EN24}
\newcommand{\DOCTITRE}{Lesson 24 — Writing: Writes compositions on taking part and managing time in community/workplace/school activities}
% OnBuch+ — préambule « cours signés » ENRICHI (compilé avec tectonic / XeLaTeX)
% Système visuel « Pop » d'OnBuch+ : fond crème (jamais blanc), bordures encre
% épaisses, ombres « dures » décalées, titres Archivo Black en capitales.
% Version MATHÉMATIQUES : graphes (pgfplots/tikz), boîtes pédagogiques
% (définition, propriété, méthode, attention, le savais-tu, exercice résolu,
% exercice, TP, bilan), page de garde et signature OnBuch+ ancrée en bas.
%
% Macros attendues dans main.tex AVANT \input{preamble} :
%   \DOCMATIERE  \DOCNIVEAU  \DOCMODULE  \DOCLECON (numéro)  \DOCTITRE
\documentclass[11pt]{article}

\usepackage{fontspec}
\usepackage[english,french]{babel} % français = langue principale ; l'anglais via \eng{..} / textebox
\usepackage[a4paper,margin=2cm,top=3.4cm,bottom=2.8cm,headheight=1.4cm,footskip=1.3cm]{geometry}
\usepackage{amsmath,amssymb,amsfonts}
\usepackage{mathtools} % \xmapsto, \coloneqq... souvent utilisés par les modèles
\usepackage{cancel} % simplifications barrées (bilans de forces)
% \abs{z} (module, valeur absolue) et \norm{u} (norme), utilisés par les modèles
\providecommand{\abs}{}\let\abs\relax\DeclarePairedDelimiter{\abs}{\lvert}{\rvert}
\providecommand{\norm}{}\let\norm\relax\DeclarePairedDelimiter{\norm}{\lVert}{\rVert}
\usepackage[table]{xcolor} % [table] : fournit aussi \rowcolors (alternance de lignes)
\usepackage{pagecolor}
\usepackage{tikz}
\usetikzlibrary{arrows.meta,positioning,calc,shapes.geometric,shapes.misc,shapes.arrows,decorations.pathmorphing,decorations.pathreplacing,decorations.markings,patterns,shadows,mindmap,trees,fit,backgrounds,matrix,intersections,angles,quotes}
\usepackage{pgfplots}
\usepackage[version=4]{mhchem} % équations nucléaires \ce{^{235}_{92}U}
\usepackage[european,siunitx]{circuitikz} % schémas électriques
\pgfplotsset{compat=1.18}
\usepgfplotslibrary{fillbetween,groupplots} % aires entre courbes (calcul intégral)
\usepackage{enumitem}
\usepackage{fancyhdr}
% [most] charge toutes les bibliothèques tcolorbox (listings, minted, raster,
% documentation...) et gonfle inutilement la mémoire TeX sur un document long
% (~300 boîtes) : on ne charge que ce qui est réellement utilisé.
\usepackage[skins,breakable]{tcolorbox}
\usepackage{graphicx}
\usepackage{colortbl}
\usepackage{booktabs}
\usepackage{tabularx}
\usepackage{diagbox} % case d'en-tête diagonale des tableaux à double entrée
% Colonnes extensibles centrée / gauche / droite (C, L, R), souvent utilisées par les modèles
\newcolumntype{C}{>{\centering\arraybackslash}X}
\newcolumntype{L}{>{\raggedright\arraybackslash}X}
\newcolumntype{R}{>{\raggedleft\arraybackslash}X}
\usepackage{array}
\usepackage{multicol}
\usepackage{multirow}
\usepackage{siunitx}
% parse-numbers=false : accepte aussi \SI{2,5 \times 10^{-3}}{...}
\sisetup{locale=FR,output-decimal-marker={,},group-minimum-digits=4,parse-numbers=false}
\DeclareSIUnit\ohm{\ensuremath{\symup{\Omega}}} % U+2126 absent de la police
\DeclareSIUnit\division{div} % réglages d'oscilloscope (ms/div, V/div)
\DeclareSIUnit\tr{tr} % tours (vitesse de rotation, ex. \SI{1500}{\tr\per\minute})
\DeclareSIUnit\molar{M} % concentration molaire (ex. \SI{3}{\molar}, \SI{1}{\micro\molar})
\DeclareSIUnit\an{an} % année (le modèle confond parfois avec \year, un compteur TeX) — ex. \SI{5}{\kilo\meter\per\an}
\DeclareSIUnit\FCFA{FCFA} % franc CFA (ex. \SI{500}{\FCFA\per\kilo\gram})
\DeclareSIUnit\degreeBrix{\textdegree Brix} % teneur en sucre d'un sirop/jus (réfractométrie)
\DeclareSIUnit\month{mois} % le modèle utilise parfois \month, un compteur TeX, comme unité de durée
\usepackage{qrcode}
% \degree : commande fréquemment utilisée par les modèles pour un angle ou une
% température en dehors de \SI{}{} (non fournie par les paquets chargés).
\providecommand{\degree}{^{\circ}}
% \qed (fin de démonstration), utilisé par les modèles sans amsthm
\providecommand{\qed}{\hfill$\square$}
% \barre : double barre des valeurs interdites dans un tableau de variations (tkz-tab)
\providecommand{\barre}{\ensuremath{\|}}
% environnement proof (amsthm non chargé) utilisé par les modèles
\ifdefined\proof\else\newenvironment{proof}[1][Démonstration]{\par\noindent\textit{#1.}\ }{\qed\par}\fi
\usepackage{lastpage}
\usepackage{hyperref}
\hypersetup{hidelinks}

% ============================================================
% Palette « Pop » OnBuch+ — mêmes hex que la classe P (plus_ui.dart)
% ============================================================
\definecolor{poporange}{HTML}{F59321}
\definecolor{popdark}{HTML}{E07A0C}
\definecolor{popcream}{HTML}{FFF6E8}
\definecolor{popink}{HTML}{1C1714}
\definecolor{poppurple}{HTML}{7A5AE0}
\definecolor{popgreen}{HTML}{2E9E6A}
\definecolor{poppink}{HTML}{E0457B}
\definecolor{popblue}{HTML}{2F7DE1}
\definecolor{popgold}{HTML}{C98A12}
\definecolor{popmuted}{HTML}{8A7F76}
\definecolor{popline}{HTML}{EADFD2}
\definecolor{popgray}{HTML}{8A7F76}
\colorlet{popgrey}{popgray}
% Alias tolérants (le modèle nomme parfois les couleurs différemment)
\colorlet{popwhite}{white}
\colorlet{popblack}{popink}
\colorlet{popred}{poppink}
\colorlet{popyellow}{popgold}
% Teintes claires (fonds de tableaux/encadrés) — le modèle nomme parfois
% les couleurs avec un suffixe L (ex. popblueL) sans les avoir définies.
\colorlet{poporangeL}{poporange!20}
\colorlet{popdarkL}{popdark!20}
\colorlet{poppurpleL}{poppurple!20}
\colorlet{popgreenL}{popgreen!20}
\colorlet{poppinkL}{poppink!20}
\colorlet{popblueL}{popblue!20}
\colorlet{popgoldL}{popgold!20}
\colorlet{popredL}{popred!20}
\colorlet{popgrayL}{popgray!20}
% Teintes claires pour fonds de tableaux / zones
\colorlet{popblueL}{popblue!10!white}
\colorlet{poporangeL}{poporange!14!white}
\colorlet{popgreenL}{popgreen!12!white}
\colorlet{poppurpleL}{poppurple!12!white}
\colorlet{poppinkL}{poppink!12!white}
\colorlet{popgoldL}{popgold!14!white}

\pagecolor{popcream}

% ============================================================
% Typographie
% ============================================================
\defaultfontfeatures{Ligatures=TeX}
\setmainfont{PlusJakartaSans-Medium.ttf}[Path=../../../fonts/,BoldFont=PlusJakartaSans-Bold.ttf]
% Maths en sans-serif (Fira Math) avec chiffres et lettres droites de Plus
% Jakarta, pour que les formules se fondent dans le texte.
\usepackage[math-style=ISO,bold-style=ISO]{unicode-math}
\setmathfont{FiraMath-Regular.otf}[Path=../../../fonts/]
\setmathfont{PlusJakartaSans-Medium.ttf}[Path=../../../fonts/,range={up/num,up/{Latin,latin},\mathup}]
\usepackage{icomma} % virgule décimale sans espace en mode math
% Caractères mathématiques Unicode absents des polices de texte (Plus Jakarta,
% Archivo Black) : rendus par leur équivalent mathématique, en texte comme en
% formule — sinon ils s'affichent en carré vide (ex. « Arithmétique dans ℤ »).
\usepackage{newunicodechar}
\newunicodechar{ℝ}{\ensuremath{\mathbb{R}}}
\newunicodechar{ℤ}{\ensuremath{\mathbb{Z}}}
\newunicodechar{ℂ}{\ensuremath{\mathbb{C}}}
\newunicodechar{ℕ}{\ensuremath{\mathbb{N}}}
\newunicodechar{ℚ}{\ensuremath{\mathbb{Q}}}
\newunicodechar{𝔻}{\ensuremath{\mathbb{D}}}
\newunicodechar{α}{\ensuremath{\alpha}}
\newunicodechar{β}{\ensuremath{\beta}}
\newunicodechar{γ}{\ensuremath{\gamma}}
\newunicodechar{δ}{\ensuremath{\delta}}
\newunicodechar{ε}{\ensuremath{\varepsilon}}
\newunicodechar{θ}{\ensuremath{\theta}}
\newunicodechar{λ}{\ensuremath{\lambda}}
\newunicodechar{μ}{\ensuremath{\mu}}
\newunicodechar{σ}{\ensuremath{\sigma}}
\newunicodechar{τ}{\ensuremath{\tau}}
\newunicodechar{φ}{\ensuremath{\varphi}}
\newunicodechar{ψ}{\ensuremath{\psi}}
\newunicodechar{ω}{\ensuremath{\omega}}
\newunicodechar{Σ}{\ensuremath{\Sigma}}
% majuscules produites par \MakeUppercase dans les titres (α-aminés -> Α)
\newunicodechar{Α}{\ensuremath{\alpha}}
\newunicodechar{Β}{\ensuremath{\beta}}
\newunicodechar{Γ}{\ensuremath{\Gamma}}
\newunicodechar{Δ}{\ensuremath{\Delta}}
\newunicodechar{Ε}{\ensuremath{\varepsilon}}
\newunicodechar{Θ}{\ensuremath{\Theta}}
\newunicodechar{Λ}{\ensuremath{\Lambda}}
\newunicodechar{Μ}{\ensuremath{\mu}}
\newunicodechar{Τ}{\ensuremath{\tau}}
\newunicodechar{Φ}{\ensuremath{\Phi}}
\newunicodechar{Ψ}{\ensuremath{\Psi}}
\newunicodechar{Ω}{\ensuremath{\Omega}}
\newunicodechar{↦}{\ensuremath{\mapsto}}
\newunicodechar{∈}{\ensuremath{\in}}
\newunicodechar{∉}{\ensuremath{\notin}}
\newunicodechar{∧}{\ensuremath{\wedge}}
\newunicodechar{⇔}{\ensuremath{\Leftrightarrow}}
\newunicodechar{⟺}{\ensuremath{\Longleftrightarrow}}
\newunicodechar{⇒}{\ensuremath{\Rightarrow}}
\newunicodechar{⟹}{\ensuremath{\Longrightarrow}}
\newunicodechar{∘}{\ensuremath{\circ}}
\newunicodechar{∪}{\ensuremath{\cup}}
\newunicodechar{∩}{\ensuremath{\cap}}
\newunicodechar{≡}{\ensuremath{\equiv}}
\newunicodechar{⊂}{\ensuremath{\subset}}
\newunicodechar{⊥}{\ensuremath{\perp}}
\newunicodechar{∃}{\ensuremath{\exists}}
\newunicodechar{∀}{\ensuremath{\forall}}
\newunicodechar{∛}{\ensuremath{\sqrt[3]{\,}}}
\newunicodechar{∜}{\ensuremath{\sqrt[4]{\,}}}
\newunicodechar{⃗}{\ensuremath{{}^{\to}}}
\newunicodechar{ⁿ}{\ifmmode^{n}\else\textsuperscript{\ensuremath{n}}\fi}
\newunicodechar{ˣ}{\ifmmode^{x}\else\textsuperscript{\ensuremath{x}}\fi}
\newunicodechar{ᵇ}{\ifmmode^{b}\else\textsuperscript{\ensuremath{b}}\fi}
\newunicodechar{ʳ}{\ifmmode^{r}\else\textsuperscript{\ensuremath{r}}\fi}
\newunicodechar{ᵘ}{\ifmmode^{u}\else\textsuperscript{\ensuremath{u}}\fi}
\newunicodechar{ʸ}{\ifmmode^{y}\else\textsuperscript{\ensuremath{y}}\fi}
\newunicodechar{ᵃ}{\ifmmode^{a}\else\textsuperscript{\ensuremath{a}}\fi}
\newunicodechar{ᵉ}{\ifmmode^{e}\else\textsuperscript{\ensuremath{e}}\fi}
\newunicodechar{ᵏ}{\ifmmode^{k}\else\textsuperscript{\ensuremath{k}}\fi}
\newunicodechar{ᵖ}{\ifmmode^{p}\else\textsuperscript{\ensuremath{p}}\fi}
\newunicodechar{ᴺ}{\ifmmode^{N}\else\textsuperscript{\ensuremath{N}}\fi}
\newunicodechar{ᶜ}{\ifmmode^{c}\else\textsuperscript{\ensuremath{c}}\fi}
\newunicodechar{ʰ}{\ifmmode^{h}\else\textsuperscript{\ensuremath{h}}\fi}
\newunicodechar{ᵠ}{\ifmmode^{q}\else\textsuperscript{\ensuremath{q}}\fi}
\newunicodechar{ₙ}{\ifmmode_{n}\else\textsubscript{\ensuremath{n}}\fi}
\newunicodechar{ₐ}{\ifmmode_{a}\else\textsubscript{\ensuremath{a}}\fi}
\newunicodechar{ᵢ}{\ifmmode_{i}\else\textsubscript{\ensuremath{i}}\fi}
\newunicodechar{ᵣ}{\ifmmode_{r}\else\textsubscript{\ensuremath{r}}\fi}
\newunicodechar{ₖ}{\ifmmode_{k}\else\textsubscript{\ensuremath{k}}\fi}
\newunicodechar{ᵦ}{\ifmmode_{\beta}\else\textsubscript{\ensuremath{\beta}}\fi}
\newunicodechar{ₓ}{\ifmmode_{x}\else\textsubscript{\ensuremath{x}}\fi}
\newfontfamily\popdisplay{ArchivoBlack-Regular.ttf}[Path=../../../fonts/]
\newfontfamily\poplogo{SpaceGrotesk-Bold.ttf}[Path=../../../fonts/]
\newfontfamily\popsemi{PlusJakartaSans-SemiBold.ttf}[Path=../../../fonts/]
\newfontfamily\popxbold{PlusJakartaSans-ExtraBold.ttf}[Path=../../../fonts/]
\newfontfamily\popxboldls{PlusJakartaSans-ExtraBold.ttf}[Path=../../../fonts/,LetterSpace=3.0]

\setlist[enumerate]{leftmargin=1.7em,itemsep=5pt,topsep=4pt}
\setlist[itemize]{leftmargin=1.7em,itemsep=4pt,topsep=4pt,label={\color{poporange}\textbullet}}
\setlength{\parindent}{0pt}
\setlength{\parskip}{5pt}
\renewcommand{\arraystretch}{1.25}

% pgfplots : style Pop par défaut
\pgfplotsset{
  every axis/.append style={axis line style={popink,line width=1pt},tick style={popink},
    grid style={popline,line width=0.5pt},label style={font=\small},tick label style={font=\footnotesize}},
  popcurve/.style={poporange,line width=1.8pt,no markers,smooth},
}
% Même style disponible dans un \draw TikZ simple (hors pgfplots)
\tikzset{popcurve/.style={poporange,line width=1.8pt}}
\tikzset{popgrid/.style={popline,very thin}}
\colorlet{popcurve}{poporange} % le nom du style est parfois utilisé comme couleur

% ============================================================
% Pill / squiggle
% ============================================================
\newcommand{\poppill}[3]{%
  \tikz[baseline=-0.55ex]\node[draw=popink,line width=1.2pt,rounded corners=2.6pt,%
    fill=#2,inner xsep=6.5pt,inner ysep=3pt]{%
    \popxboldls\fontsize{7}{7}\selectfont\color{#3}\MakeUppercase{#1}};%
}
\newcommand{\popsquiggle}[1]{%
  \tikz[baseline=-0.4ex]{
    \draw[#1,line width=2.6pt,line cap=round]
      plot [smooth] coordinates {(0,0.09) (0.34,0.01) (0.68,0.21) (1.02,0.01) (1.36,0.10)};
  }%
}
% Mot-clé surligné (marqueur orange sous le texte)
\newcommand{\cle}[1]{{\popxbold\color{popdark}#1}}

% ============================================================
% En-tête / pied
% ============================================================
\pagestyle{fancy}
\fancyhf{}
\renewcommand{\headrulewidth}{0pt}
\renewcommand{\footrulewidth}{0pt}
\lhead{\raisebox{-0.15ex}{\poplogo\fontsize{13}{13}\selectfont\color{popink}OnBuch\color{poporange}+}}
\rhead{\poppill{\DOCMATIERE\ \textbullet\ \DOCNIVEAU\ \textbullet\ Leçon \DOCLECON}{white}{popink}}
\lfoot{\popxbold\fontsize{7.6}{7.6}\selectfont\color{popmuted}\MakeUppercase{Cours signé OnBuch+}\ \textbullet\ {\popsemi\DOCTITRE}}
\rfoot{\tikz[baseline=-0.6ex]\node[rounded corners=8.5pt,fill=popink,minimum height=17pt,minimum width=17pt,inner xsep=5pt,inner sep=0]{\popxbold\fontsize{8}{8}\selectfont\color{popcream}\thepage\,/\,\pageref*{LastPage}};}

% ============================================================
% Titres
% ============================================================
\newcommand{\chaptitre}[2]{%
  \begin{tcolorbox}[enhanced,colback=poporange,colframe=popink,boxrule=2.6pt,arc=11pt,
      drop shadow=popink,left=15pt,right=15pt,top=12pt,bottom=11pt,before skip=3pt,after skip=18pt]
    \poppill{#2}{popink}{popcream}\par\vspace{9pt}
    {\raggedright\hyphenpenalty=10000\exhyphenpenalty=10000\popdisplay\fontsize{22}{24}\selectfont\color{popcream}\MakeUppercase{#1}\par}
  \end{tcolorbox}
}
% Section numérotée : pastille orange avec le numéro + titre Archivo
\newcounter{popsec}
\newcommand{\coursec}[1]{%
  \stepcounter{popsec}\setcounter{popsub}{0}%
  \par\vspace{16pt}\noindent
  \begin{minipage}{\linewidth}
  \tikz[baseline=-0.7ex]\node[draw=popink,line width=1.6pt,fill=poporange,rounded corners=4pt,minimum size=22pt,inner sep=0]{\popdisplay\fontsize{12}{12}\selectfont\color{popcream}\thepopsec};%
  \hspace{8pt}{\popdisplay\fontsize{15}{17}\selectfont\color{popink}\MakeUppercase{#1}}\par
  \vspace{4pt}\noindent{\color{poporange}\rule{36pt}{3.6pt}}
  \end{minipage}\par\nopagebreak\vspace{8pt}
}
\newcounter{popsub}
\newcommand{\courssub}[1]{%
  \stepcounter{popsub}%
  \par\vspace{9pt}\noindent{\popxbold\fontsize{11.5}{13}\selectfont\color{popink}{\color{poporange}\thepopsec.\thepopsub}\hspace{6pt}#1}\par\nopagebreak\vspace{4pt}
}

% ============================================================
% Boîtes pédagogiques — même recette : fond blanc, bordure épaisse colorée,
% ombre dure encre, badge plein. Titre optionnel [#1].
% ============================================================
\tcbset{popbase/.style={breakable,enhanced,colback=white,boxrule=2.3pt,arc=9pt,
  shadow={3pt}{-3pt}{0pt}{popink},left=14pt,right=14pt,top=11pt,bottom=11pt,before skip=11pt,after skip=15pt,
  fontupper=\popsemi\fontsize{10.6}{14.5}\selectfont\color{popink},
  fontlower=\popsemi\fontsize{10.6}{14.5}\selectfont\color{popink}}}
% Coche dessinée (le glyphe ✓ n'existe pas dans les polices du document)
\newcommand{\popcheck}[1]{\tikz[baseline=-0.5ex]\draw[#1,line width=1.6pt,line cap=round,line join=round](-0.13,0.0)--(-0.04,-0.09)--(0.13,0.1);}
\providecommand{\coche}{\popcheck{popgreen}}
\providecommand{\resultat}[1]{\par\smallskip\noindent\fcolorbox{popgreen}{popgreenL}{\parbox{\dimexpr\linewidth-2\fboxsep-2\fboxrule\relax}{\textbf{Résultat.} #1}}\par\smallskip}
\newcommand{\popboxhead}[3]{\poppill{#1}{#2}{white}\ifx\relax#3\relax\else\hspace{6pt}{\popxbold\fontsize{10.6}{12}\selectfont\color{popink}#3}\fi\par\vspace{7pt}}

\newtcolorbox{aretenir}[1][]{popbase,colframe=popgreen,before upper={\popboxhead{À retenir}{popgreen}{#1}}}
% Texte en anglais (document de lecture, dialogue, extrait) : cadre
% d'étude. Un titre optionnel (source, auteur) s'affiche dans l'en-tête.
\newtcolorbox{textebox}[1][]{popbase,colframe=popblue,before upper={\popboxhead{Text}{popblue}{#1}\begin{otherlanguage}{english}},after upper={\end{otherlanguage}}}
% Marques ✓ / ✗ (forme correcte / fautive) souvent utilisées par les modèles
\providecommand{\cross}{\textcolor{poppink}{\ensuremath{\times}}}
\providecommand{\xmark}{\cross}\providecommand{\crossmark}{\cross}\providecommand{\cmark}{\ensuremath{\checkmark}}
% Ligne de réponse / justification à compléter (inventée par les modèles)
\providecommand{\justification}[1]{\par\noindent\textit{Justification~:} \dotfill\par}
% \textipa (package tipa non chargé) : la police ne couvre pas l'API -> simple passage
\providecommand{\textipa}[1]{#1}
% Soulignement (ulem non chargé) : \uline, \ul
\providecommand{\uline}[1]{\underline{#1}}\providecommand{\ul}[1]{\underline{#1}}
% \glossaire{\item ...} inventée par les modèles : liste de glossaire
\providecommand{\glossaire}[1]{\begin{itemize}#1\end{itemize}}
% \alert (beamer) inventée par les modèles : texte mis en évidence
\providecommand{\alert}[1]{\textbf{\textcolor{poppink}{#1}}}
% variantes ASCII d'ordinaux écrites en macro
\providecommand{\iere}{\textsuperscript{re}}\providecommand{\ieme}{\textsuperscript{e}}\providecommand{\ere}{\textsuperscript{re}}\providecommand{\eme}{\textsuperscript{e}}\providecommand{\er}{\textsuperscript{er}}\providecommand{\ie}{\textsuperscript{e}}
% Ordinaux français écrits en macro par les modèles : 1\ière, 3\ième
\newcommand{\ière}{\textsuperscript{re}}\newcommand{\ième}{\textsuperscript{e}}
% \exemple (inventée par les modèles) : étiquette « Exemple : »
\providecommand{\exemple}{\textbf{Exemple~:}\ }
% \square utilisable aussi hors mode math (cases à cocher)
\AtBeginDocument{\let\squareMath\square\renewcommand{\square}{\ensuremath{\squareMath}}}
% Mot ou phrase en anglais à l'intérieur d'un paragraphe français
\newcommand{\eng}[1]{\ifmmode\text{\textit{#1}}\else\begingroup\selectlanguage{english}\itshape #1\endgroup\fi}
\let\esp\eng % alias toléré
\newtcolorbox{exemplebox}[1][]{popbase,colframe=poporange,before upper={\popboxhead{Exemple}{poporange}{#1}}}
\newtcolorbox{definition}[1][]{popbase,colframe=popblue,before upper={\popboxhead{Définition}{popblue}{#1}}}
\newtcolorbox{propriete}[1][]{popbase,colframe=poppurple,before upper={\popboxhead{Propriété}{poppurple}{#1}}}
\newtcolorbox{methode}[1][]{popbase,colframe=popgold,before upper={\popboxhead{Méthode}{popgold}{#1}}}
\newtcolorbox{attention}[1][]{popbase,colframe=poppink,before upper={\popboxhead{Attention}{poppink}{#1}}}
\newtcolorbox{experience}[1][]{popbase,colframe=popblue,colback=popblueL,before upper={\popboxhead{Expérience}{popblue}{#1}}}
% « Le savais-tu ? » : carte violette pleine (contexte, culture, Cameroun)
\newtcolorbox{savaistu}[1][]{popbase,colback=poppurple,colframe=popink,
  fontupper=\popsemi\fontsize{10.4}{14.2}\selectfont\color{white},
  before upper={\poppill{Le savais-tu\,?}{white}{poppurple}\ifx\relax#1\relax\else\hspace{6pt}{\popxbold\color{white}#1}\fi\par\vspace{7pt}}}
% Exercice résolu : énoncé en haut, \tcblower puis la solution sur fond crème
\newcounter{popexo}\newcounter{popexr}
\newtcolorbox{exoresolu}[1][]{popbase,colframe=poporange,colbacklower=popcream,
  segmentation style={popink,line width=1.2pt,dashed},
  before upper={\stepcounter{popexr}\popboxhead{Exercice résolu \thepopexr}{poporange}{#1}},
  before lower={\poppill{Solution}{popink}{popcream}\par\vspace{6pt}}}
% Exercice (non résolu en place) — la correction est dans la partie Corrigés
\newtcolorbox{exercice}[1][]{popbase,colframe=popink,
  before upper={\stepcounter{popexo}\popboxhead{Exercice \thepopexo}{popink}{#1}}}
% Corrigé
\newtcolorbox{corrige}[1][]{popbase,colframe=popgreen,colback=popgreenL,
  before upper={\popboxhead{Corrigé}{popgreen}{#1}}}
% Objectifs / prérequis / bilan
\newtcolorbox{objectifs}{popbase,colframe=popink,colback=poporangeL,
  before upper={\popboxhead{Objectifs de la leçon}{popink}{}}}
\newtcolorbox{prerequis}{popbase,colframe=popink,colback=popblueL,
  before upper={\popboxhead{Prérequis}{popblue}{}}}
\newtcolorbox{bilan}{popbase,colframe=popink,colback=white,boxrule=2.8pt,
  before upper={\popboxhead{Fiche bilan}{poporange}{L'essentiel à connaître pour l'examen}}}
% Figure encadrée avec légende
\newtcolorbox{popfigure}[1][]{enhanced,breakable=false,colback=white,colframe=popline,boxrule=1.4pt,arc=8pt,
  left=10pt,right=10pt,top=10pt,bottom=8pt,before skip=10pt,after skip=12pt,halign=center,
  lower separated=false,halign lower=center,
  fontlower=\popsemi\fontsize{9}{11.5}\selectfont\color{popmuted},
  #1}
\newcommand{\legende}[1]{\par\vspace{6pt}{\popsemi\fontsize{9}{11.5}\selectfont\color{popmuted}#1}}

% ============================================================
% Signature OnBuch+
% ============================================================
\newcommand{\poplogoinline}[1]{{\poplogo\fontsize{#1}{#1}\selectfont\color{popink}OnBuch\color{poporange}+}}
\newcommand{\signatureonbuch}{%
  \begin{tcolorbox}[enhanced,colback=popink,colframe=popink,boxrule=2.2pt,arc=10pt,
      drop shadow=poporange,left=14pt,right=14pt,top=10pt,bottom=10pt,before skip=0pt,after skip=0pt]
    \tikz[baseline=-0.7ex]\node[circle,fill=popgreen,draw=popcream,line width=1.3pt,minimum size=22pt,inner sep=0]{\popcheck{white}};%
    \hspace{9pt}\begin{minipage}[c]{0.62\linewidth}
      {\popxbold\fontsize{12}{13}\selectfont\color{popcream}Signé\ \textbullet\ {\poplogo OnBuch\color{poporange}+}}\par\vspace{2pt}
      {\raggedright\popsemi\fontsize{8}{10.5}\selectfont\color{popline}\MakeUppercase{Équipe pédagogique OnBuch+}\\\MakeUppercase{Conforme au programme officiel MINESEC}\par}
    \end{minipage}\hfill
    {\poplogo\fontsize{22}{22}\selectfont\color{popcream}OnBuch\color{poporange}+}
  \end{tcolorbox}%
}

% ============================================================
% Page de garde
% #1 titre · #2 matière · #3 niveau · #4 module · #5 n° leçon · #6 durée ·
% #7 description / objectifs (texte ou itemize)
% ============================================================
\newcommand{\pagegarde}[7]{%
  \thispagestyle{empty}
  \newgeometry{a4paper,margin=1.7cm,top=1.6cm,bottom=1.4cm}%
  \begin{tikzpicture}[remember picture,overlay]
    \fill[poporange!18] ([xshift=-3.2cm,yshift=-3.6cm]current page.north east) circle (4.6cm);
    \fill[poppurple!14] ([xshift=2.4cm,yshift=7.6cm]current page.south west) circle (3.8cm);
  \end{tikzpicture}%
  {\poplogo\fontsize{30}{30}\selectfont\color{popink}OnBuch\color{poporange}+}\hfill
  \raisebox{0.6ex}{\poppill{Cours signé \textbullet\ Premium}{popink}{popcream}}\par
  \vspace{1pt}\popsquiggle{poppurple}\par
  \vspace{8pt}
  \begin{tcolorbox}[enhanced,colback=poporange,colframe=popink,boxrule=2.8pt,arc=13pt,
      drop shadow=popink,left=18pt,right=18pt,top=12pt,bottom=13pt,after skip=10pt]
    \poppill{#2 \textbullet\ #3}{popink}{popcream}\hspace{5pt}\poppill{#4}{white}{popink}\par\vspace{8pt}
    {\popdisplay\fontsize{14}{14}\selectfont\color{popink}LEÇON #5}\par\vspace{4pt}
    {\raggedright\hyphenpenalty=10000\exhyphenpenalty=10000\popdisplay\fontsize{25}{27}\selectfont\color{popcream}\MakeUppercase{#1}\par}
    \vspace{8pt}
    {\popxbold\fontsize{9.5}{9.5}\selectfont\color{popink}\MakeUppercase{Durée conseillée : #6}}
  \end{tcolorbox}
  \begin{tcolorbox}[enhanced,colback=white,colframe=popink,boxrule=2.2pt,arc=10pt,
      drop shadow=popink,left=16pt,right=16pt,top=10pt,bottom=10pt,after skip=8pt]
    \poppill{Dans cette leçon}{poppurple}{white}\par\vspace{5pt}
    \popsemi\fontsize{9.3}{12.6}\selectfont\color{popink}#7
  \end{tcolorbox}
  \noindent\begin{minipage}[t]{0.487\linewidth}
    \begin{tcolorbox}[enhanced,colback=popgreen,colframe=popink,boxrule=2.2pt,arc=10pt,
        drop shadow=popink,left=11pt,right=11pt,top=10pt,bottom=10pt,height=3.1cm,valign=center]
      \tcbox[colback=white,colframe=popink,boxrule=1.3pt,arc=5pt,boxsep=0pt,nobeforeafter,
        left=3pt,right=3pt,top=3pt,bottom=3pt,valign=center]{\qrcode[height=1.9cm]{https://play.google.com/store/apps/details?id=com.luvvix.onbuch&hl=fr}}%
      \hspace{8pt}\begin{minipage}[c]{0.5\linewidth}
        {\popdisplay\fontsize{11}{12.5}\selectfont\color{white}\MakeUppercase{Télécharge OnBuch}}\par\vspace{4pt}
        {\raggedright\popsemi\fontsize{8.4}{11}\selectfont\color{white}Gratuit \textbullet\ Google Play\\Tuteur IA, annales, concours, résultats\par}
      \end{minipage}
    \end{tcolorbox}
  \end{minipage}\hfill
  \begin{minipage}[t]{0.487\linewidth}
    \begin{tcolorbox}[enhanced,colback=poppurple,colframe=popink,boxrule=2.2pt,arc=10pt,
        drop shadow=popink,left=11pt,right=11pt,top=10pt,bottom=10pt,height=3.1cm,valign=center]
      \tikz[baseline=-0.7ex]\node[circle,fill=white,draw=popink,line width=1.3pt,minimum size=1.6cm,inner sep=0]{\popdisplay\fontsize{24}{24}\selectfont\color{poppurple}+};%
      \hspace{8pt}\begin{minipage}[c]{0.6\linewidth}
        {\popdisplay\fontsize{11}{12.5}\selectfont\color{white}\MakeUppercase{Passe à OnBuch+}}\par\vspace{4pt}
        {\raggedright\popsemi\fontsize{8.4}{11}\selectfont\color{white}Tous les cours signés, Tuteur IA illimité, fiches PDF — onglet OnBuch+ de l'app\par}
      \end{minipage}
    \end{tcolorbox}
  \end{minipage}
  \vspace{8pt}
  \popcontenu
  \vfill
  \signatureonbuch
  \restoregeometry
  \newpage
}

% Rangée « Ce cours contient » (page de garde)
\newcommand{\poptile}[3]{% #1 couleur #2 gros texte #3 légende
  \begin{tcolorbox}[enhanced,colback=#1,colframe=popink,boxrule=2pt,arc=9pt,shadow={2.5pt}{-2.5pt}{0pt}{popink},
      left=6pt,right=6pt,top=9pt,bottom=9pt,halign=center,valign=center,height=1.75cm,width=\linewidth,nobeforeafter]
    {\popdisplay\fontsize{12.5}{13}\selectfont\color{white}\MakeUppercase{#2}}\par\vspace{3pt}
    {\centering\popsemi\fontsize{7.8}{9.5}\selectfont\color{white}#3\par}
  \end{tcolorbox}}
\newcommand{\popcontenu}{%
  \par\noindent{\popxboldls\fontsize{7.5}{7.5}\selectfont\color{popmuted}CE COURS SIGNÉ CONTIENT}\par\vspace{7pt}
  \noindent\begin{minipage}[t]{0.235\linewidth}\poptile{popblue}{Cours}{complet, illustré, pas à pas}\end{minipage}\hfill
  \begin{minipage}[t]{0.235\linewidth}\poptile{poporange}{Méthodes}{et exercices résolus}\end{minipage}\hfill
  \begin{minipage}[t]{0.235\linewidth}\poptile{poppink}{Exercices}{type examen + corrigés}\end{minipage}\hfill
  \begin{minipage}[t]{0.235\linewidth}\poptile{popgreen}{Fiche bilan}{l'essentiel à retenir}\end{minipage}\par}

% Sommaire
\newtcolorbox{sommaire}{enhanced,colback=white,colframe=popink,boxrule=2.2pt,arc=10pt,
  drop shadow=popink,left=18pt,right=18pt,top=13pt,bottom=13pt,before skip=6pt,after skip=20pt,
  before upper={\poppill{Sommaire}{poppurple}{white}\par\vspace{9pt}\popsemi\fontsize{10.6}{14.5}\selectfont\color{popink}}}

% Signature de fin de cours — poussée tout en bas de la dernière page
\newcommand{\findecours}{%
  \par\vfill\nopagebreak
  \begin{center}\popsquiggle{poporange}\end{center}\vspace{4pt}
  \signatureonbuch
}
\AtBeginDocument{\def\llbracket{\mathopen{[\mkern-3.5mu[}}\def\rrbracket{\mathclose{]\mkern-3.5mu]}}} % intervalles d'entiers (glyphe ⟦ absent de la police)
% Glyphes absents de Fira Math : redéfinis à partir de glyphes disponibles
\AtBeginDocument{%
  \def\setminus{\mathbin{\backslash}}\let\smallsetminus\setminus
  \def\vdots{\mathord{\vbox{\baselineskip3.2pt\lineskiplimit0pt\kern4pt\hbox{.}\hbox{.}\hbox{.}}}}%
  \def\ddots{\mathinner{\mkern1mu\raise6.5pt\hbox{.}\mkern2mu\raise3.6pt\hbox{.}\mkern2mu\raise0.7pt\hbox{.}\mkern1mu}}%
  \def\triangle{\mathord{\scalebox{0.9}{$\symup{\Delta}$}}}%
  \def\nabla{\mathord{\raisebox{\depth}{\scalebox{1}[-1]{$\symup{\Delta}$}}}}%
  \def\checkmark{\popcheck{popdark}}%
  \def\star{\mathbin{\ast}}\def\bigstar{\mathord{\ast}}%
  \def\diamond{\mathbin{\scalebox{0.6}{$\Diamond$}}}%
  \def\bigcup{\mathop{\mathchoice{\raisebox{-0.3em}{\scalebox{1.7}{$\cup$}}}{\scalebox{1.25}{$\cup$}}{\cup}{\cup}}\displaylimits}%
  \def\bigcap{\mathop{\mathchoice{\raisebox{-0.3em}{\scalebox{1.7}{$\cap$}}}{\scalebox{1.25}{$\cap$}}{\cap}{\cap}}\displaylimits}%
  \def\lesseqqgtr{\mathrel{\lessgtr}}\def\bigoplus{\mathop{\oplus}\displaylimits}%
}
\providecommand{\ssi}{si et seulement si} % abréviation fréquente
\AtBeginDocument{\providecommand{\wideparen}{\overparen}} % arc de cercle

%% FIGFIX-BEGIN v4 — correctifs globaux des figures (docs/onbuch-generation/tools/figfix)
\usepackage{adjustbox}\usepackage{etoolbox}
% toute figure TikZ/pgfplots plus large que la ligne est réduite à la largeur disponible
\BeforeBeginEnvironment{tikzpicture}{\begin{adjustbox}{max width=\linewidth}}
\AfterEndEnvironment{tikzpicture}{\end{adjustbox}}
% légendes pgfplots lisibles par-dessus les courbes
\pgfplotsset{every axis/.append style={legend style={fill=white,fill opacity=0.92,text opacity=1,draw=black!15,inner sep=3pt}}}
% étiquettes d'arêtes (syntaxe edge["..."]) avec fond blanc
\tikzset{every edge quotes/.append style={fill=white,inner sep=1.2pt}}
%% FIGFIX-END
\begin{document}

\pagegarde{Lesson 24 — Writing: Writes compositions on taking part and managing time in community/workplace/school activities}{Anglais}{Tle A}{Chapter 2 : Economic Life and Occupations}{EN24}{2 h}{Cette leçon d'expression écrite apprend aux élèves de Tle A à rédiger une composition claire et bien structurée sur la participation et la gestion du temps dans les activités communautaires, scolaires ou professionnelles. Elle fournit la structure, les connecteurs, le vocabulaire, un modèle commenté et une grille d'évaluation conformes aux attentes du Bac.
\vspace{4pt}
\begin{multicols}{2}\small
\begin{itemize}
\item À la fin de cette leçon, je sais analyser un sujet de composition et dégager les consignes essentielles.
\item Je sais construire un plan en trois parties : introduction, développement, conclusion.
\item Je sais rédiger une introduction avec une accroche, une présentation du sujet et une annonce de plan.
\item Je sais développer des paragraphes articulés par des connecteurs logiques variés.
\item Je sais utiliser le vocabulaire de la participation et de la gestion du temps.
\item Je sais employer les temps verbaux appropriés (present simple, past simple, futur) dans une composition.
\end{itemize}\end{multicols}}

\begin{sommaire}
\begin{multicols}{2}
\begin{enumerate}
\item Analyse du sujet et planification
\item La structure d'une composition
\item Connecteurs logiques et temps verbaux
\item Vocabulaire utile : participation et gestion du temps
\item Modèle commenté et erreurs à éviter
\item Grille d'évaluation et entraînement au Bac
\item Activité d'intégration
\item Méthodes et exercices résolus
\item Exercices
\item Corrigés des exercices
\item Fiche bilan
\end{enumerate}
\end{multicols}
\end{sommaire}

\begin{objectifs}
\begin{itemize}
\item À la fin de cette leçon, je sais analyser un sujet de composition et dégager les consignes essentielles.
\item Je sais construire un plan en trois parties : introduction, développement, conclusion.
\item Je sais rédiger une introduction avec une accroche, une présentation du sujet et une annonce de plan.
\item Je sais développer des paragraphes articulés par des connecteurs logiques variés.
\item Je sais utiliser le vocabulaire de la participation et de la gestion du temps.
\item Je sais employer les temps verbaux appropriés (present simple, past simple, futur) dans une composition.
\item Je sais rédiger une conclusion qui résume et ouvre le débat.
\item Je sais relire et corriger ma production à l'aide d'une grille d'évaluation.
\item Je sais éviter les erreurs fréquentes des candidats francophones au Bac.
\end{itemize}
\end{objectifs}

\begin{prerequis}
\begin{itemize}
\item Les temps de base : present simple, past simple, be going to / will (Première).
\item Les connecteurs simples : and, but, because, so, first, then, finally (Seconde/Première).
\item Le vocabulaire de la vie scolaire et des activités quotidiennes (Première).
\item La rédaction d'un paragraphe simple avec phrase d'introduction (topic sentence).
\item La ponctuation et la majuscule en anglais.
\end{itemize}
\end{prerequis}

\newpage

\coursec{Analyse du sujet et planification}

\begin{textebox}[Your first task as class prefect]
Imagine: your school in Bamenda is organising a “Clean-Up Day” next Saturday, and the principal asks you, as class prefect, to write a short composition explaining how students will take part and how the day will be timed. You sit at your desk, pen in hand, and two questions come to your mind: what exactly must I write, and how do I organise my ideas?

In the UK or the USA, students write similar essays about volunteering and time management for their final exams. At the Baccalauréat, you are often asked to write a composition on community, school or workplace activities. Knowing how to plan, organise and correct such an essay is the key to top marks.

Many candidates lose marks not because their English is weak, but because they start writing immediately, without a plan. Their ideas are good, but they are mixed up. The examiner cannot follow the argument, and the mark falls. In this lesson, we learn the first and most important step of writing: \emph{analysing the topic and planning the composition}.
\end{textebox}

\textbf{Problématique.} Comment passer d'un sujet d'examen à un plan clair, en quelques minutes, avant même d'écrire la première phrase ? C'est toute la question de cette première section : \emph{lire}, \emph{comprendre}, \emph{brainstormer}, \emph{sélectionner}, \emph{planifier}.

\courssub{Lire le sujet : la méthode WHO–WHAT–HOW–HOW LONG}

\begin{definition}[Topic, task, text type]
En anglais, le \cle{topic} est le thème du sujet (ici : \eng{community, school or workplace activities}), la \cle{task} est ce qu'on vous demande de faire (\eng{write about taking part and managing time}), et le \cle{text type} est le genre de texte attendu (ici : a \cle{composition}, c'est-à-dire un essai scolaire organisé en paragraphes).
\end{definition}

La première règle d'or : \textbf{lisez le sujet deux fois}. À la première lecture, vous comprenez le sens général. À la deuxième lecture, vous \textbf{soulignez les mots-clés}. Prenons un sujet typique du Bac :

\begin{exemplebox}[Un sujet typique à analyser]
\eng{Write a composition of about 250 words on the following topic: “Taking part in school or community activities and managing your time well are the keys to success.” Do you agree? Illustrate your answer with examples from your school or your neighbourhood.}

Mots-clés à souligner :
\begin{itemize}
\item \eng{composition} → type de texte : essai en paragraphes, pas une lettre ;
\item \eng{about 250 words} → longueur attendue ;
\item \eng{taking part} et \eng{managing your time} → les deux idées obligatoires ;
\item \eng{Do you agree?} → il faut donner une opinion ;
\item \eng{examples from your school or your neighbourhood} → il faut des exemples concrets (Bamenda, Douala, votre lycée).
\end{itemize}
\end{exemplebox}

\begin{methode}[Analyser un sujet en quatre questions : WHO – WHAT – HOW – HOW LONG]
Avant d'écrire, répondez par écrit à quatre questions :
\begin{enumerate}
\item \textbf{WHO} — Pour qui j'écris ? (\eng{Who am I writing for?}) Ici : l'examinateur, ou le principal et les camarades. Le destinataire détermine le ton.
\item \textbf{WHAT} — Quel est le thème exact ? (\eng{What is the topic?}) Ici : la participation (\eng{taking part}) et la gestion du temps (\eng{managing time}) dans les activités scolaires, communautaires ou professionnelles.
\item \textbf{HOW} — Quel type de texte, quel ton ? (\eng{What text type, what tone?}) Ici : une composition, au ton \textbf{neutre ou semi-formel} : pas d'abréviations (\eng{don't} est toléré mais préférez \eng{do not}), pas d'argot, phrases complètes.
\item \textbf{HOW LONG} — Combien de mots ? (\eng{How many words?}) Respectez la fourchette indiquée : pour \eng{about 250 words}, écrivez entre 220 et 280 mots environ.
\end{enumerate}
\end{methode}

\begin{attention}[Le registre : ni trop familier, ni trop solennel]
Une composition scolaire exige un ton \cle{semi-formel}. Les francophones font deux erreurs opposées :
\begin{itemize}
\item trop familier : \eng{Hi guys, cleaning is cool!} — à éviter ;
\item maladroitement traduit du français soutenu : \eng{I have the honour to inform you that...} — réservé aux lettres officielles.
\end{itemize}
Écrivez simplement : \eng{In my opinion, taking part in community activities is very important for students.} (« À mon avis, participer aux activités communautaires est très important pour les élèves. »)
\end{attention}

\courssub{Le brainstorming : produire des idées en vrac}

Une fois le sujet compris, ne rédigez pas encore. Pendant deux ou trois minutes, notez \textbf{6 à 8 idées en vrac}, en anglais, sans les juger. C'est le \cle{brainstorming} (prononcé « brène-stor-ming »), littéralement « tempête de cerveau ».

\begin{exemplebox}[Brainstorming sur notre sujet]
\begin{itemize}
\item \eng{clean-up campaign in the neighbourhood} (campagne de nettoyage du quartier)
\item \eng{study group before the exam} (groupe d'étude avant l'examen)
\item \eng{sports day at school} (journée sportive)
\item \eng{part-time job at the market} (petit travail au marché)
\item \eng{making a timetable} (établir un emploi du temps)
\item \eng{setting priorities} (fixer des priorités)
\item \eng{learning responsibility and teamwork} (apprendre la responsabilité et le travail d'équipe)
\item \eng{balancing studies and activities} (équilibrer études et activités)
\end{itemize}
\end{exemplebox}

\begin{exemplebox}[Deux phrases-clés à retenir]
\eng{Taking part in community activities teaches young people responsibility.} « Participer aux activités communautaires apprend la responsabilité aux jeunes. »

\eng{Managing my time well helps me succeed at school.} « Bien gérer mon temps m'aide à réussir à l'école. »

Remarquez la construction : le sujet de la phrase est un \textbf{groupe verbal en -ing} (\eng{taking part}, \eng{managing my time}), ce qui correspond au français « participer », « bien gérer ». En anglais, l'infinitif ne peut pas être sujet : on dit \eng{Taking part is important}, jamais \emph{To take part is important} dans ce contexte courant.
\end{exemplebox}

\courssub{Sélectionner et classer les idées}

Vous ne pouvez pas tout dire en 250 mots. Il faut \textbf{sélectionner 3 idées principales} et les \textbf{classer du plus simple au plus important}. Cette progression donne de la force à votre argumentation : le lecteur monte en puissance avec vous.

\begin{center}
\begin{tabularx}{\linewidth}{|X|X|X|}
\hline
\rowcolor{popblueL} Étape & English & Français \\
\hline
Idée 1 (la plus simple) & \eng{taking part in a clean-up campaign} & participer à une campagne de nettoyage \\
\hline
Idée 2 (plus large) & \eng{making a timetable and setting priorities} & établir un emploi du temps et des priorités \\
\hline
Idée 3 (la plus importante) & \eng{learning responsibility for future work life} & apprendre la responsabilité pour la vie professionnelle future \\
\hline
\end{tabularx}
\end{center}

\begin{attention}[Faux amis et pièges de vocabulaire]
\begin{itemize}
\item \eng{to take part in} = participer à. Ne dites jamais \emph{to participate to} : le verbe \eng{participate} se construit avec \textbf{in} : \eng{I participate in the sports day.}
\item \eng{to manage my time} = gérer mon temps. Attention : \eng{to manage to do something} signifie « réussir à faire quelque chose » : \eng{I managed to finish on time.} (« J'ai réussi à finir à temps. »)
\item \eng{a timetable} = un emploi du temps ; \eng{a schedule} = un programme, un planning (plus américain). Ne confondez pas avec \eng{a watch} (une montre) !
\item \eng{actually} = en fait, réellement ; « actuellement » se dit \eng{currently} ou \eng{at present}.
\end{itemize}
\end{attention}

\courssub{Construire le plan en trois parties}

Toute composition au Bac suit la même architecture : \textbf{introduction / développement / conclusion}. Le développement contient 2 ou 3 paragraphes, et chaque paragraphe s'ouvre sur une \cle{topic sentence}, une « phrase-guide » qui annonce l'idée du paragraphe.

\begin{definition}[Topic sentence]
La \cle{topic sentence} (phrase-guide) est la première phrase d'un paragraphe. Elle annonce l'idée principale du paragraphe. Toutes les autres phrases du paragraphe la développent ou l'illustrent. Exemple : \eng{First of all, taking part in community activities teaches students responsibility.} — tout le paragraphe parlera ensuite de la responsabilité.
\end{definition}

\begin{popfigure}
\begin{tikzpicture}[mindmap, grow cyclic, every node/.style={concept, align=center, font=\small},
root concept/.append style={concept color=popblue, text=white, font=\small\bfseries},
level 1/.append style={level distance=3.4cm, sibling angle=90},
level 1 concept/.append style={concept color=poporange, text=popdark}]
\node[root concept]{My composition}
  child { node[align=center]{Introduction\\ hook + opinion} }
  child { node[align=center]{Body 1:\\ taking part\\ (clean-up day)} }
  child { node[align=center]{Body 2:\\ managing time\\ (timetable)} }
  child { node[align=center]{Conclusion\\ summary + opening} };
\end{tikzpicture}
\legende{Carte mentale du plan : un centre, quatre branches, une idée par paragraphe.}
\end{popfigure}

\begin{exemplebox}[Un plan au brouillon, avec une topic sentence par paragraphe]
\textbf{Introduction} — \eng{In my school in Bamenda, students take part in many activities, and I believe that participation and good time management are the keys to success.}

\textbf{Body 1} — Topic sentence : \eng{First of all, taking part in community activities, such as our Clean-Up Day, teaches young people responsibility and teamwork.}

\textbf{Body 2} — Topic sentence : \eng{Secondly, managing time well, with a clear timetable and priorities, helps students succeed in their studies.}

\textbf{Conclusion} — \eng{To conclude, students who take part in activities and manage their time become responsible adults, ready for work life.}
\end{exemplebox}

\courssub{Gérer son temps à l'examen}

La gestion du temps s'applique aussi à vous, candidat ! Pour une épreuve de 50 minutes consacrée à la composition, voici la répartition recommandée.

\begin{popfigure}
\begin{tikzpicture}[font=\small]
\draw[-{Stealth}, thick, popdark] (0,0) -- (12.6,0);
\foreach \x/\t/\c in {0/{Start}/popdark, 1.2/{5 min}/popblue, 3.6/{15 min}/poporange, 10.8/{45 min}/popgreen, 12/{50 min}/poppurple} {
  \fill[\c] (\x,0) circle (0.09);
  \node[below, popdark] at (\x,-0.1) {\t};
}
\draw[thick, popblue] (0,0.35) -- (1.2,0.35) node[midway, above, popblue]{analyse};
\draw[thick, poporange] (1.2,0.35) -- (3.6,0.35) node[midway, above, poporange]{plan};
\draw[thick, popgreen] (3.6,0.35) -- (10.8,0.35) node[midway, above, popgreen]{rédaction};
\draw[thick, poppurple] (10.8,0.35) -- (12,0.35) node[midway, above, poppurple]{relecture};
\end{tikzpicture}
\legende{Frise du temps de travail : 5 min d'analyse, 10 min de plan, 30 min de rédaction, 5 min de relecture.}
\end{popfigure}

\begin{aretenir}[Les cinq réflexes du bon candidat]
\begin{enumerate}
\item Lire le sujet \textbf{deux fois} et souligner les mots-clés.
\item Répondre aux quatre questions : WHO – WHAT – HOW – HOW LONG.
\item Brainstormer 6 à 8 idées, puis n'en garder que 3, classées du plus simple au plus important.
\item Écrire un plan au brouillon avec une topic sentence par paragraphe.
\item Garder 5 minutes pour la relecture finale.
\end{enumerate}
\end{aretenir}

\begin{attention}[Ne rédigez jamais sans plan]
Au Bac, une composition désorganisée perd des points \textbf{même avec un bon anglais}. L'examinateur note la pertinence, l'organisation et la correction de la langue. Un candidat qui écrit directement mélange ses idées, se répète, oublie la conclusion — et dépasse souvent le temps. Dix minutes de plan vous font gagner des points, pas en perdre.
\end{attention}

\begin{savaistu}[La règle des « 5 paragraphs »]
Au Royaume-Uni, les élèves qui préparent le GCSE (l'équivalent du brevet et du début du lycée) apprennent la règle des \eng{five paragraphs} : une introduction, trois paragraphes de développement, une conclusion. Aux États-Unis, on parle du \eng{five-paragraph essay}. Cette structure est très proche de celle attendue au Baccalauréat camerounais : les bonnes habitudes d'écriture sont les mêmes dans tout le monde anglophone !
\end{savaistu}

\begin{exoresolu}[Analyser un sujet de Bac]
Énoncé — Voici un sujet : \eng{Write a composition of about 250 words: “Students should take part in community activities and learn to manage their time.” Discuss, giving examples from your town.} 1) Soulignez les mots-clés. 2) Répondez aux quatre questions WHO – WHAT – HOW – HOW LONG. 3) Proposez trois idées classées.
\tcblower
Solution détaillée —
1) Mots-clés : \eng{composition}, \eng{about 250 words}, \eng{take part in community activities}, \eng{manage their time}, \eng{discuss}, \eng{examples from your town}.

2) WHO : l'examinateur du Bac. WHAT : la participation aux activités communautaires et la gestion du temps. HOW : une composition au ton semi-formel, avec une opinion personnelle (\eng{discuss} = discuter, donner des arguments). HOW LONG : environ 250 mots (entre 220 et 280).

3) Trois idées classées du plus simple au plus important : (a) \eng{taking part in a clean-up campaign in Douala} — exemple concret et simple ; (b) \eng{using a timetable to balance studies and activities} — la gestion du temps ; (c) \eng{becoming responsible adults prepared for the workplace} — l'idée la plus large, qui ouvre sur l'avenir professionnel. Ce classement donne une progression logique : du geste quotidien à la préparation de la vie adulte.
\end{exoresolu}

\begin{exercice}[À vous de planifier]
Sujet : \eng{Write a composition of about 250 words on: “Taking part in school activities and managing time well prepare students for the world of work.”}
\begin{enumerate}
\item Soulignez les mots-clés du sujet et répondez aux quatre questions WHO – WHAT – HOW – HOW LONG.
\item Faites un brainstorming de 6 idées en anglais.
\item Sélectionnez 3 idées et classez-les du plus simple au plus important.
\item Rédigez le plan au brouillon : une topic sentence pour l'introduction, une pour chaque paragraphe du développement, une pour la conclusion.
\end{enumerate}
\end{exercice}

\begin{corrige}[Exercice : À vous de planifier]
1) Mots-clés : \eng{composition}, \eng{250 words}, \eng{taking part in school activities}, \eng{managing time}, \eng{world of work}. WHO : l'examinateur. WHAT : participation et gestion du temps comme préparation au monde du travail. HOW : composition semi-formelle. HOW LONG : environ 250 mots.

2) Brainstorming possible : \eng{sports day, debate club, clean-up campaign, study group, making a timetable, setting priorities, teamwork, learning punctuality}.

3) Trois idées classées : (a) \eng{taking part in the debate club} ; (b) \eng{using a timetable to balance studies and activities} ; (c) \eng{learning teamwork and punctuality for the workplace}.

4) Plan possible — Introduction : \eng{In my opinion, school activities and good time management are excellent preparation for the world of work.} Body 1 : \eng{First of all, taking part in activities such as the debate club builds confidence.} Body 2 : \eng{Secondly, managing time with a timetable teaches students to respect deadlines.} Conclusion : \eng{To conclude, active and well-organised students become reliable workers.}
\end{corrige}

\coursec{La structure d'une composition}

Dans la section précédente, nous avons appris à analyser le sujet et à planifier nos idées avant d'écrire. Vous savez maintenant quelles idées vous voulez défendre. Reste une question essentielle : comment les organiser sur la copie ? Une composition anglaise, comme un bâtiment bien construit, repose sur une architecture précise et immuable : une \cle{introduction}, un \cle{développement} (\eng{body}) et une \cle{conclusion}. C'est cette architecture que nous allons étudier en détail, car au Baccalauréat, la clarté de la structure compte autant que la correction de la langue.

\courssub{Vue d'ensemble : les trois parties d'une composition}

\begin{definition}[The structure of a composition]
Une \cle{composition} anglaise (\eng{composition} ou \eng{essay}) est un texte argumentatif ou narratif organisé en trois parties :
\begin{itemize}
\item the \cle{introduction} : elle présente le sujet et annonce le plan ;
\item the \cle{body} (développement) : il développe les idées, un paragraphe par idée ;
\item the \cle{conclusion} : elle résume les idées et donne une opinion personnelle.
\end{itemize}
\end{definition}

\begin{popfigure}
\begin{center}
\begin{tikzpicture}[node distance=0.35cm]
\node[draw=popblue, fill=popblueL, rounded corners, text width=11.5cm, align=center, minimum height=1.6cm] (intro) {\textbf{\textcolor{popdark}{INTRODUCTION}}\\[2pt] \eng{Hook} (accroche) + \eng{topic} (sujet) + \eng{plan} (annonce du plan)};
\node[draw=poporange, fill=poporangeL, rounded corners, text width=11.5cm, align=center, minimum height=2.6cm, below=of intro] (body) {\textbf{\textcolor{popdark}{BODY — DEVELOPMENT}}\\[2pt] \eng{Paragraph 1} : topic sentence + explanation + example\\ \eng{Paragraph 2} : topic sentence + explanation + example\\ \eng{Paragraph 3} : topic sentence + explanation + example};
\node[draw=popgreen, fill=popgreenL, rounded corners, text width=11.5cm, align=center, minimum height=1.6cm, below=of body] (concl) {\textbf{\textcolor{popdark}{CONCLUSION}}\\[2pt] \eng{Summary} (résumé) + \eng{personal opinion} (opinion) + \eng{opening} (ouverture)};
\draw[-{Stealth}, thick, popblue] (intro.south) -- (body.north);
\draw[-{Stealth}, thick, poporange] (body.south) -- (concl.north);
\end{tikzpicture}
\end{center}
\legende{L'architecture d'une composition : trois blocs empilés, chacun avec sa fonction.}
\end{popfigure}

\begin{attention}[Ne pas écrire les titres des parties]
Contrairement à une dissertation française, on n'écrit \textbf{jamais} les mots \eng{Introduction}, \eng{Development} ou \eng{Conclusion} comme titres dans la copie. Les parties se reconnaissent par leur \textbf{contenu} et par les \textbf{alinéas} (ou sauts de ligne) qui séparent les paragraphes. Écrire ces titres est une erreur qui coûte des points au Bac.
\end{attention}

\begin{savaistu}[Pourquoi cette structure est-elle si stricte ?]
Dans la tradition scolaire anglo-saxonne, on apprend très tôt le modèle du \eng{five-paragraph essay} (dissertation en cinq paragraphes) : une introduction, trois paragraphes de développement, une conclusion. L'idée est que le lecteur doit pouvoir suivre le raisonnement sans effort : il sait toujours où il en est. Au Baccalauréat camerounais, les correcteurs appliquent exactement cette grille de lecture : une copie bien structurée est immédiatement repérée et valorisée.
\end{savaistu}

\courssub{L'introduction : trois à quatre phrases}

\begin{methode}[Observer avant d'écrire]
Lisons d'abord une introduction modèle, puis décomposons-la phrase par phrase. C'est en observant des modèles réussis que l'on apprend à écrire, pas en inventant sa propre méthode le jour de l'examen.
\end{methode}

\begin{textebox}[Model introduction — read and analyse]
Have you ever wondered why some students succeed in everything they do? The secret is simple: they take part in school activities and manage their time well. This composition will show how participation and time management lead to success.
\end{textebox}

Observons : la première phrase est une \textbf{question} qui capte l'attention du lecteur ; la deuxième \textbf{présente le sujet} (participation et gestion du temps) ; la troisième \textbf{annonce le plan} (\eng{This composition will show...}). Trois phrases, trois fonctions. C'est exactement ce que l'on attend de vous.

\begin{propriete}[Les trois éléments de l'introduction]
Une introduction efficace contient, dans l'ordre :
\begin{enumerate}
\item \textbf{The hook} (l'accroche) : une phrase qui attire l'attention du lecteur ;
\item \textbf{The topic} (la présentation du sujet) : on reformule le sujet avec ses propres mots ;
\item \textbf{The plan} (l'annonce du plan) : une phrase qui annonce ce que la composition va démontrer.
\end{enumerate}
Longueur : \textbf{3 à 4 phrases} maximum. Une introduction trop longue déséquilibre toute la copie.
\end{propriete}

\textbf{Les types d'accroches.} Il existe plusieurs façons d'accrocher le lecteur. En voici les trois plus sûres au Bac :

\begin{center}
\begin{tabularx}{\linewidth}{|l|X|X|}
\hline
\rowcolor{popblueL} \textbf{Type d'accroche} & \textbf{Exemple en anglais} & \textbf{Traduction} \\
\hline
Question & \eng{Have you ever taken part in a community project?} & Avez-vous déjà participé à un projet communautaire ? \\
\hline
Fait général & \eng{In many Cameroonian schools, students organise clean-up days.} & Dans de nombreuses écoles camerounaises, les élèves organisent des journées de nettoyage. \\
\hline
Citation ou proverbe & \eng{As the saying goes, “time is money”.} & Comme le dit le proverbe, « le temps, c'est de l'argent ». \\
\hline
\end{tabularx}
\end{center}

\begin{exemplebox}[Annoncer le plan : formules utiles]
\begin{itemize}
\item \eng{This composition will show how participation and time management lead to success.} — « Cette composition montrera comment la participation et la gestion du temps mènent au succès. »
\item \eng{In this essay, I will explain why taking part in community activities is important.} — « Dans cette dissertation, j'expliquerai pourquoi participer aux activités communautaires est important. »
\item \eng{I will first discuss the benefits of participation, and then explain how to manage one's time.} — « Je traiterai d'abord des avantages de la participation, puis j'expliquerai comment gérer son temps. »
\end{itemize}
\end{exemplebox}

\begin{attention}[Pièges de l'introduction]
\begin{itemize}
\item Ne recopiez \textbf{jamais} le sujet mot pour mot : reformulez-le. Le correcteur veut voir que vous l'avez compris.
\item Évitez \eng{I am going to write a composition about...} : trop scolaire. Préférez \eng{This composition will show...}.
\item N'annoncez pas un plan que vous ne tiendrez pas : le développement doit suivre exactement l'ordre annoncé.
\item Attention à l'accroche-proverbe : n'utilisez que des proverbes dont vous êtes \textbf{certain} de la formulation anglaise. Un proverbe traduit mot à mot du français (\eng{“who can the more can the less”}) est pire qu'une simple question.
\end{itemize}
\end{attention}

\courssub{Le développement : un paragraphe, une idée}

\begin{definition}[The topic sentence]
La \cle{topic sentence} (phrase d'idée principale) est la \textbf{première phrase d'un paragraphe}. Elle annonce l'idée principale du paragraphe ; toutes les autres phrases du paragraphe la développent, l'expliquent ou l'illustrent. Si une phrase ne se rattache pas à la topic sentence, elle doit être supprimée ou déplacée.
\end{definition}

\begin{propriete}[La structure d'un paragraphe type]
Chaque paragraphe du développement suit le schéma en quatre temps :
\begin{enumerate}
\item \textbf{Topic sentence} : l'idée principale, introduite par un connecteur (\eng{First of all}, \eng{Moreover}, \eng{Finally}) ;
\item \textbf{Explanation} : deux ou trois phrases qui expliquent l'idée ;
\item \textbf{Example} : un exemple concret, souvent introduit par \eng{For example} ou \eng{For instance} ;
\item \textbf{Mini-conclusion} (facultative mais recommandée) : une phrase qui boucle le paragraphe et prépare la transition.
\end{enumerate}
Règle d'or : \textbf{un paragraphe = une seule idée}. Au Bac, deux ou trois paragraphes de développement suffisent.
\end{propriete}

\begin{exemplebox}[Un paragraphe complet, décomposé]
\eng{First of all, taking part in school activities builds team spirit.} (topic sentence) \eng{When students work together in a club or on a project, they learn to listen to others and to share responsibilities.} (explanation) \eng{For example, members of the football team at Government High School Buea train together every Friday and learn to trust one another.} (example) \eng{In this way, school activities prepare students for life in society.} (mini-conclusion)

Traduction : « Tout d'abord, participer aux activités scolaires développe l'esprit d'équipe. Quand les élèves travaillent ensemble dans un club ou sur un projet, ils apprennent à écouter les autres et à partager les responsabilités. Par exemple, les membres de l'équipe de football du lycée de Buea s'entraînent ensemble chaque vendredi et apprennent à se faire confiance. Ainsi, les activités scolaires préparent les élèves à la vie en société. »
\end{exemplebox}

\begin{attention}[La cohérence : chaque phrase s'accroche à la précédente]
Une erreur fréquente des francophones est de juxtaposer des phrases sans lien logique. En anglais, chaque phrase doit se rattacher clairement à la précédente, grâce à un connecteur (\eng{Moreover}, \eng{For example}, \eng{As a result}) ou à une reprise (\eng{this habit}, \eng{these activities}). Relisez chaque paragraphe en vous demandant : « Pourquoi cette phrase est-elle ici ? » Si vous ne pouvez pas répondre, supprimez-la.
\end{attention}

\courssub{La conclusion : deux à trois phrases}

\begin{propriete}[Les trois fonctions de la conclusion]
La conclusion est la partie la plus courte (2 à 3 phrases). Elle contient :
\begin{enumerate}
\item un \textbf{résumé} des idées principales, introduit par \eng{In conclusion} ou \eng{To sum up} ;
\item une \textbf{opinion personnelle} : \eng{In my opinion...}, \eng{I believe that...} ;
\item éventuellement une \textbf{ouverture} : une question ou une perspective plus large.
\end{enumerate}
Interdiction absolue : n'introduisez \textbf{aucune idée nouvelle} dans la conclusion.
\end{propriete}

\begin{exemplebox}[Exemples de conclusions traduits]
\eng{In conclusion, managing time is the key to success.} — « En conclusion, gérer son temps est la clé du succès. »

\eng{To sum up, participation in community activities teaches young people responsibility and solidarity. In my opinion, every student should take part in at least one activity.} — « En résumé, la participation aux activités communautaires apprend aux jeunes la responsabilité et la solidarité. À mon avis, chaque élève devrait participer à au moins une activité. »
\end{exemplebox}

\begin{attention}[Faux amis et formules à éviter en conclusion]
\begin{itemize}
\item \eng{In my opinion} est correct ; ne dites jamais \eng{according to me} (calque de « selon moi »), qui est fautif en anglais standard. Dites \eng{in my view} ou \eng{I think that...}.
\item N'écrivez pas \eng{Finally} pour ouvrir la conclusion : \eng{finally} marque le dernier point d'une énumération. Pour conclure, utilisez \eng{In conclusion}, \eng{To sum up} ou \eng{All in all}.
\end{itemize}
\end{attention}

\courssub{Longueur et présentation matérielle}

\begin{aretenir}[Les chiffres à connaître pour le Bac]
\begin{itemize}
\item Longueur totale : \textbf{150 à 250 mots} selon la consigne (respectez toujours la fourchette indiquée) ;
\item Introduction : 3 à 4 phrases ; développement : 2 à 3 paragraphes ; conclusion : 2 à 3 phrases ;
\item Les paragraphes doivent être \textbf{visibles} : alinéa (retrait de première ligne) ou saut de ligne entre les paragraphes ;
\item Comptez vos mots approximativement (comptez une ligne, multipliez par le nombre de lignes).
\end{itemize}
\end{aretenir}

\begin{center}
\begin{tabularx}{\linewidth}{|l|X|X|}
\hline
\rowcolor{popblueL} \textbf{Partie} & \textbf{Contenu} & \textbf{Longueur indicative} \\
\hline
Introduction & hook + topic + plan & 3–4 phrases (environ 40 mots) \\
\hline
Body (P1, P2, P3) & topic sentence + explanation + example & 4–6 phrases par paragraphe \\
\hline
Conclusion & summary + opinion (+ opening) & 2–3 phrases (environ 30 mots) \\
\hline
\end{tabularx}
\end{center}

\courssub{Entraînement guidé}

\begin{exoresolu}[Identifier les éléments d'une introduction]
Énoncé — Voici une introduction. Identifiez l'accroche, la présentation du sujet et l'annonce du plan : \eng{In many villages around Bamenda, young people spend their free time helping their community. Taking part in community activities is both useful and rewarding. This essay will explain the benefits of participation and show how to manage one's time.}
\tcblower
Solution — Phrase 1 : \eng{In many villages around Bamenda...} est l'\textbf{accroche} (fait général). Phrase 2 : \eng{Taking part in community activities is both useful and rewarding} \textbf{présente le sujet}. Phrase 3 : \eng{This essay will explain... and show...} \textbf{annonce le plan} en deux parties, que le développement devra suivre dans le même ordre.
\end{exoresolu}

\begin{exercice}[À vous de jouer]
Sujet : \eng{Write a composition on how students can take part in school activities and still manage their time well.}
\begin{enumerate}
\item Rédigez une accroche sous forme de question.
\item Rédigez la phrase de présentation du sujet.
\item Rédigez l'annonce du plan en deux parties.
\item Rédigez la topic sentence du premier paragraphe de développement.
\end{enumerate}
\end{exercice}

\begin{corrige}[Exercice — éléments de réponse]
\begin{enumerate}
\item \eng{Have you ever wondered how some students do everything and still get good marks?}
\item \eng{Taking part in school activities is possible if students learn to manage their time.}
\item \eng{This composition will first show the benefits of participation, and then explain how to organise one's time.}
\item \eng{First of all, school activities help students to develop useful skills.}
\end{enumerate}
\end{corrige}

Maintenant que l'architecture de la composition est claire, il nous reste à la faire vivre : la section suivante présentera les connecteurs logiques et les temps verbaux qui relient ces blocs entre eux.

\coursec{Connecteurs logiques et temps verbaux}

Dans la section précédente, nous avons appris à construire le squelette d'une composition : une introduction qui présente le sujet, un développement en deux ou trois paragraphes, et une conclusion. Mais un squelette sans articulations ne bouge pas. Ce sont les \cle{connecteurs logiques} (\eng{linking words}) qui relient vos idées entre elles et qui donnent à votre texte son rythme et sa clarté. À cela s'ajoute le choix rigoureux des \cle{temps verbaux}, qui situe chaque idée dans le temps. C'est précisément ce que les correcteurs du Baccalauréat récompensent sous le critère \eng{coherence and cohesion}.

\courssub{Observer avant de comprendre : deux versions d'un même texte}

Lisez attentivement les deux versions suivantes. Les idées sont exactement les mêmes ; seuls les connecteurs et les temps changent.

\begin{textebox}[Version A — without linkers]
I joined the environmental club last year. I was busy with my studies. I found time to help. We cleaned the school yard every Friday. We planted flowers near the classrooms. The school became cleaner. I learned to manage my time. I am going to organise a recycling campaign next term.
\end{textebox}

\begin{textebox}[Version B — with linkers and correct tenses]
First of all, I joined the environmental club last year. Although I was busy with my studies, I found time to help. Every Friday, we cleaned the school yard; moreover, we planted flowers near the classrooms. As a result, the school became cleaner. This experience taught me that time management helps students succeed. Therefore, I am going to organise a recycling campaign next term.
\end{textebox}

\begin{definition}[Glossaire]
\begin{itemize}
\item \eng{to join} (verbe) : rejoindre, s'inscrire à
\item \eng{although} (conjonction) : bien que
\item \eng{moreover} (adverbe) : de plus, en outre
\item \eng{as a result} (locution) : par conséquent
\item \eng{time management} (nom) : gestion du temps
\item \eng{therefore} (adverbe) : donc, c'est pourquoi
\item \eng{a recycling campaign} (nom) : une campagne de recyclage
\end{itemize}
\end{definition}

La version B est plus fluide, plus mature, plus « baccalauréat ». Analysons maintenant chaque famille de connecteurs.

\courssub{Les connecteurs d'addition}

Ils servent à ajouter une idée à une idée précédente, comme le français « de plus », « en outre », « également ».

\begin{propriete}[Connecteurs d'addition : forme et emploi]
\begin{itemize}
\item \eng{first of all} : pour introduire le premier argument. \eng{First of all, taking part in school activities builds confidence.} « Tout d'abord, participer aux activités scolaires développe la confiance en soi. »
\item \eng{moreover} et \eng{furthermore} : pour ajouter un argument plus fort (registre soutenu, idéal au Bac). Suivis d'une virgule en début de phrase. \eng{Moreover, volunteering develops leadership skills.} « De plus, le bénévolat développe les qualités de leader. »
\item \eng{in addition} : synonyme de \eng{moreover}, légèrement moins formel. \eng{In addition, we organised a football match.} « En outre, nous avons organisé un match de football. »
\item \eng{also} : plus simple ; se place généralement au milieu de la phrase, avant le verbe principal. \eng{The club also helps new students.} « Le club aide aussi les nouveaux élèves. »
\end{itemize}
\end{propriete}

\begin{attention}[Also : attention à la place]
En français, « aussi » peut commencer une phrase (« Aussi avons-nous décidé... »). En anglais, \eng{also} ne commence presque jamais une phrase dans une rédaction scolaire. Écrivez \eng{We also decided to...} et non « Also we decided to... ». Pour débuter une phrase, utilisez \eng{moreover} ou \eng{in addition}.
\end{attention}

\courssub{Les connecteurs d'ordre et de temps}

Ils organisent la chronologie d'un récit ou d'une description d'activités : essentiels pour raconter une journée de participation (\eng{clean-up day}, \eng{Sports Day}, réunion de club).

\begin{propriete}[Connecteurs d'ordre et de temps]
\begin{itemize}
\item \eng{first} : d'abord. \eng{First, we met in the school hall.} « D'abord, nous nous sommes retrouvés dans le hall. »
\item \eng{then} : puis, ensuite. \eng{Then, the headmaster gave a short speech.} « Puis, le proviseur a fait un petit discours. »
\item \eng{next} : ensuite (après une étape). \eng{Next, we divided into teams.} « Ensuite, nous nous sommes répartis en équipes. »
\item \eng{after that} : après cela. \eng{After that, we started cleaning the classrooms.} « Après cela, nous avons commencé à nettoyer les salles. »
\item \eng{meanwhile} : pendant ce temps. \eng{Meanwhile, other students were painting the gate.} « Pendant ce temps, d'autres élèves peignaient le portail. »
\item \eng{finally} : finalement, enfin. \eng{Finally, we shared a meal together.} « Finalement, nous avons partagé un repas. »
\end{itemize}
\end{propriete}

\begin{attention}[Firstly : correct mais lourd]
Beaucoup de francophones écrivent \eng{firstly} par calque de « premièrement ». Le mot existe, mais il sonne lourd et scolaire en anglais moderne. Préférez \eng{first} ou \eng{first of all}. Évitez aussi la série mécanique \eng{firstly, secondly, thirdly} : les correcteurs la trouvent artificielle.
\end{attention}

\courssub{Les connecteurs de cause et de conséquence}

Ils expriment le « pourquoi » et le « donc » : indispensables pour argumenter sur l'importance de la gestion du temps.

\begin{propriete}[Cause et conséquence : forme et emploi]
\textbf{Pour exprimer la cause} (suivis d'une proposition avec sujet + verbe) :
\begin{itemize}
\item \eng{because} : parce que. \eng{I joined the choir because I love singing.} « J'ai rejoint la chorale parce que j'adore chanter. »
\item \eng{since} et \eng{as} : puisque, comme (plus soutenus). \eng{Since time is limited, we must plan our week.} « Puisque le temps est limité, nous devons planifier notre semaine. »
\end{itemize}
\textbf{Pour exprimer la conséquence} (en début de phrase, suivis d'une virgule, sauf \eng{so}) :
\begin{itemize}
\item \eng{so} : donc (au milieu de la phrase, sans virgule après). \eng{The hall was full, so we met outside.} « La salle était pleine, donc nous nous sommes réunis dehors. »
\item \eng{therefore} : donc, par conséquent (soutenu). \eng{Time is precious; therefore, we must not waste it.} « Le temps est précieux ; donc, nous ne devons pas le gaspiller. »
\item \eng{as a result} : par conséquent. \eng{As a result, our school became cleaner.} « Par conséquent, notre école est devenue plus propre. »
\item \eng{that is why} : c'est pourquoi. \eng{I plan my day every morning; that is why I never miss deadlines.} « Je planifie ma journée chaque matin ; c'est pourquoi je ne manque jamais les échéances. »
\end{itemize}
\end{propriete}

\begin{attention}[Because of / because : ne confondez pas]
\eng{because} + phrase complète : \eng{I stayed at home because it was raining.} \eng{because of} + nom : \eng{I stayed at home because of the rain.} Ne dites jamais « because of it was raining ». Autre piège : en anglais, on n'écrit jamais « Because I was tired. So I went to bed. » en deux phrases : reliez-les (\eng{Because I was tired, I went to bed.}).
\end{attention}

\courssub{Les connecteurs d'opposition}

Ils introduisent un contraste, une concession ou une nuance — la marque d'un raisonnement mature.

\begin{propriete}[Connecteurs d'opposition : forme et emploi]
\begin{itemize}
\item \eng{but} : mais (au milieu de la phrase). \eng{I was tired, but I finished my work.} « J'étais fatigué, mais j'ai fini mon travail. »
\item \eng{however} : cependant (début de phrase + virgule). \eng{The task was difficult. However, we succeeded.} « La tâche était difficile. Cependant, nous avons réussi. »
\item \eng{whereas} : tandis que (compare deux réalités). \eng{Some students waste time, whereas others plan carefully.} « Certains élèves perdent leur temps, tandis que d'autres planifient soigneusement. »
\item \eng{although} : bien que (+ sujet + verbe ; jamais suivi de \eng{but} dans la même phrase). \eng{Although I was busy, I found time to help.} « Bien que j'étais occupé, j'ai trouvé le temps d'aider. »
\item \eng{on the other hand} : d'un autre côté. \eng{Sports take time; on the other hand, they keep us healthy.} « Le sport prend du temps ; d'un autre côté, il nous garde en bonne santé. »
\end{itemize}
\end{propriete}

\begin{attention}[Although... but : la faute classique]
En français on entend parfois « bien que... mais ». En anglais, c'est impossible : \eng{although} et \eng{but} ne cohabitent jamais. Écrivez \eng{Although I was busy, I helped.} ou \eng{I was busy, but I helped.} — jamais « Although I was busy, but I helped. »
\end{attention}

\courssub{Les connecteurs d'exemple et de conclusion}

\begin{propriete}[Illustrer et conclure]
\begin{itemize}
\item \eng{for example} et \eng{for instance} : par exemple (phrase complète, virgule après). \eng{Many activities build team spirit. For example, the football club teaches cooperation.} « De nombreuses activités développent l'esprit d'équipe. Par exemple, le club de football enseigne la coopération. »
\item \eng{such as} : tel que (suivi d'un nom, pas d'une phrase). \eng{Clubs such as the debate club improve speaking skills.} « Des clubs tels que le club de débat améliorent l'expression orale. »
\item \eng{in conclusion} : en conclusion. \eng{In conclusion, every student should join at least one activity.} « En conclusion, chaque élève devrait participer à au moins une activité. »
\item \eng{to sum up} : pour résumer. \eng{To sum up, good time management leads to success.} « Pour résumer, une bonne gestion du temps mène à la réussite. »
\item \eng{all in all} : somme toute, tout compte fait. \eng{All in all, the experience was very rewarding.} « Tout compte fait, l'expérience fut très enrichissante. »
\end{itemize}
\end{propriete}

\begin{center}
\begin{tabularx}{\linewidth}{|l|X|X|}
\hline
\rowcolor{popblueL} \textbf{Fonction} & \textbf{Connecteurs} & \textbf{Équivalent français} \\
\hline
Addition & moreover, furthermore, in addition, also & de plus, en outre, également \\
\hline
Ordre / temps & first, then, next, after that, meanwhile, finally & d'abord, puis, ensuite, pendant ce temps, enfin \\
\hline
Cause & because, since, as & parce que, puisque, comme \\
\hline
Conséquence & so, therefore, as a result, that is why & donc, par conséquent, c'est pourquoi \\
\hline
Opposition & but, however, whereas, although, on the other hand & mais, cependant, tandis que, bien que \\
\hline
Exemple & for example, for instance, such as & par exemple, tel que \\
\hline
Conclusion & in conclusion, to sum up, all in all & en conclusion, pour résumer, somme toute \\
\hline
\end{tabularx}
\end{center}

\begin{popfigure}
\begin{tikzpicture}[font=\small]
\fill[popblueL, rounded corners] (0,0) rectangle (4.2,1.0);
\node[align=center, text=popdark] at (2.1,0.5) {\textbf{ADDITION}\\ moreover, furthermore,\\ in addition, also};
\fill[poporangeL, rounded corners] (4.8,0) rectangle (9.0,1.0);
\node[align=center, text=popdark] at (6.9,0.5) {\textbf{OPPOSITION}\\ but, however, whereas,\\ although, on the other hand};
\fill[popgreenL, rounded corners] (9.6,0) rectangle (13.4,1.0);
\node[align=center, text=popdark] at (11.5,0.5) {\textbf{CAUSE / CONSÉQUENCE}\\ because, since, so,\\ therefore, as a result};
\fill[poppurpleL, rounded corners] (2.4,-1.6) rectangle (6.6,-0.6);
\node[align=center, text=popdark] at (4.5,-1.1) {\textbf{EXEMPLE}\\ for example, for instance, such as};
\fill[popgoldL, rounded corners] (7.2,-1.6) rectangle (11.4,-0.6);
\node[align=center, text=popdark] at (9.3,-1.1) {\textbf{CONCLUSION}\\ in conclusion, to sum up, all in all};
\end{tikzpicture}
\legende{Carte mentale : chaque fonction logique a ses connecteurs.}
\end{popfigure}

\courssub{Les temps verbaux de la composition}

Une composition sur la participation et la gestion du temps mobilise trois temps principaux. Le choix du temps dépend de ce que vous exprimez : une vérité générale, une expérience passée, ou un projet.

\begin{propriete}[La règle des trois temps]
\begin{itemize}
\item \textbf{Present simple} pour les faits généraux, les vérités et les habitudes : \eng{Time management helps students.} « La gestion du temps aide les élèves. » \eng{Students who manage their time succeed.} « Les élèves qui gèrent leur temps réussissent. »
\item \textbf{Past simple} pour les événements terminés, datés, racontés : \eng{I took part in a clean-up day last month.} « J'ai participé à une journée de nettoyage le mois dernier. » \eng{Last year, I joined the school choir.} « L'année dernière, je me suis inscrit à la chorale. »
\item \textbf{be going to} pour un projet déjà décidé, une intention : \eng{I am going to join the debate club.} « Je vais rejoindre le club de débat. » \eng{Next term, I am going to organise my revision timetable.} « Le trimestre prochain, je vais organiser mon emploi du temps de révision. » \eng{will} s'emploie pour une décision spontanée ou une prédiction : \eng{I will never waste my weekends again.} « Je ne gaspillerai plus jamais mes week-ends. »
\end{itemize}
\end{propriete}

\begin{center}
\begin{tabular}{|l|l|l|}
\hline
\rowcolor{popblueL} \textbf{Temps} & \textbf{Emploi} & \textbf{Marqueurs typiques} \\
\hline
present simple & vérités générales, habitudes & every day, usually, always \\
\hline
past simple & expérience terminée & last year, yesterday, in 2024 \\
\hline
be going to & projet décidé & next term, soon, this weekend \\
\hline
will & décision spontanée, prédiction & I promise, I think... \\
\hline
\end{tabular}
\end{center}

\begin{popfigure}
\begin{tikzpicture}[font=\small]
\draw[-{Stealth}, thick, popdark] (0,0) -- (13,0);
\fill[poporange] (2.5,0) circle (0.12);
\node[align=center, text=popdark, above] at (2.5,0.25) {\textbf{past simple}\\ expérience passée\\ \emph{last year, I joined...}};
\fill[popblue] (6.5,0) circle (0.12);
\node[align=center, text=popdark, above] at (6.5,0.25) {\textbf{present simple}\\ vérités générales\\ \emph{time management helps...}};
\fill[popgreen] (10.5,0) circle (0.12);
\node[align=center, text=popdark, above] at (10.5,0.25) {\textbf{be going to}\\ projets\\ \emph{I am going to organise...}};
\node[below, text=popmuted] at (2.5,-0.15) {passé};
\node[below, text=popmuted] at (6.5,-0.15) {présent};
\node[below, text=popmuted] at (10.5,-0.15) {futur};
\end{tikzpicture}
\legende{Frise des temps : situez chaque idée de votre composition sur la ligne du temps.}
\end{popfigure}

\begin{propriete}[La concordance des temps]
Ne mélangez pas passé et présent dans un même récit sans raison. Si vous racontez une journée d'activités au passé, restez au passé du début à la fin du paragraphe : \eng{We arrived at eight, we cleaned the yard, and we planted flowers.} Le retour au présent n'est légitime que pour tirer une leçon générale de l'expérience : \eng{That day taught me that teamwork changes everything.}
\end{propriete}

\begin{attention}[Ponctuation des connecteurs]
En début de phrase, \eng{However,} \eng{Therefore,} \eng{Moreover,} \eng{For example,} et \eng{In conclusion,} sont \textbf{toujours suivis d'une virgule}. Écrivez : \eng{However, the work was not finished.} et non « However the work was not finished. » Cette virgule est un réflexe que le correcteur remarque immédiatement.
\end{attention}

\begin{methode}[Varier les connecteurs : la règle des dix]
\begin{enumerate}
\item Avant de rédiger, écrivez au brouillon une liste de \textbf{dix connecteurs} différents, couvrant toutes les fonctions : addition, ordre, cause, conséquence, opposition, exemple, conclusion.
\item Rédigez votre composition en puisant dans cette liste.
\item \textbf{Interdiction absolue} de répéter deux fois le même connecteur dans la composition : si vous avez déjà écrit \eng{moreover}, utilisez \eng{furthermore} ou \eng{in addition} la fois suivante.
\item À la relecture, entourez chaque connecteur et vérifiez : fonction correcte ? virgule après les connecteurs de début de phrase ? temps verbal cohérent ?
\end{enumerate}
\end{methode}

\begin{exoresolu}[Complétez avec le connecteur et le temps qui conviennent]
Énoncé — Complétez chaque phrase avec le connecteur proposé entre parenthèses et mettez le verbe au bon temps :
1. \eng{Last month, our class (to organise) a clean-up day. (as a result), the school yard (to become) much cleaner.}
2. \eng{(although) I (to be) very busy, I (to find) time to help my group.}
3. \eng{Time management (to help) students succeed. (therefore), I (to organise) my revision timetable next term.}
\tcblower
Solution détaillée :
1. \eng{Last month, our class \textbf{organised} a clean-up day. \textbf{As a result}, the school yard \textbf{became} much cleaner.} — \eng{last month} impose le past simple (\eng{organised}, \eng{became} ; verbe irrégulier \eng{become — became — become}). \eng{As a result} en début de phrase, suivi d'une virgule.
2. \eng{\textbf{Although} I \textbf{was} very busy, I \textbf{found} time to help my group.} — récit au passé : \eng{was}, puis \eng{found} (verbe irrégulier \eng{find — found — found}). Jamais de \eng{but} après \eng{although}.
3. \eng{Time management \textbf{helps} students succeed. \textbf{Therefore}, I \textbf{am going to organise} my revision timetable next term.} — vérité générale au present simple (\eng{helps}, avec le \eng{-s} de la troisième personne) ; projet décidé avec \eng{next term} : \eng{be going to} + base verbale.
\end{exoresolu}

\begin{exercice}[À vous de jouer]
Réécrivez le paragraphe suivant en ajoutant au moins cinq connecteurs différents et en corrigeant les temps : \eng{I join the Red Cross club last year. I am busy. I find time to help. We visit the hospital every month. We bring food to patients. I learn a lot. Next year, I train new members.}
\end{exercice}

\begin{corrige}[Exercice « À vous de jouer »]
Proposition de correction : \eng{First of all, I \textbf{joined} the Red Cross club last year. \textbf{Although} I \textbf{was} busy, I \textbf{found} time to help. Every month, we \textbf{visited} the hospital; \textbf{moreover}, we \textbf{brought} food to patients. \textbf{As a result}, I \textbf{learned} a lot about solidarity. \textbf{Therefore}, next year, I \textbf{am going to train} new members.} — Tous les événements passés sont au past simple (\eng{joined, was, found, visited, brought, learned}) ; la vérité et le projet sont respectivement au present simple et à \eng{be going to} ; cinq connecteurs variés, chacun suivi de la virgule réglementaire en début de phrase.
\end{corrige}

\begin{aretenir}[L'essentiel de la section]
\begin{itemize}
\item Chaque fonction logique a ses connecteurs : addition (\eng{moreover}), ordre (\eng{then}), cause (\eng{because}), conséquence (\eng{as a result}), opposition (\eng{however}), exemple (\eng{for instance}), conclusion (\eng{to sum up}).
\item Trois temps pour trois intentions : present simple (vérités), past simple (expérience), \eng{be going to} (projets décidés).
\item Pas de mélange passé/présent sans raison ; virgule obligatoire après les connecteurs en début de phrase ; jamais \eng{although} + \eng{but}.
\item Variez : dix connecteurs préparés au brouillon, aucun répété.
\end{itemize}
\end{aretenir}

\coursec{Vocabulaire utile : participation et gestion du temps}

Dans la section précédente, nous avons appris à relier nos idées grâce aux \cle{connecteurs logiques} et à choisir les bons temps verbaux. Mais des connecteurs sans vocabulaire précis, c'est une belle route sans voiture. Pour écrire une composition convaincante sur la participation aux activités et la gestion du temps, il faut maintenant disposer d'un \cle{champ lexical} riche, exact et bien choisi. C'est l'objectif de cette section : observer les mots en situation, les classer, apprendre leurs collocations et éviter les faux amis qui coûtent cher au Bac.

\courssub{Observer d'abord : le vocabulaire en contexte}

Lisons d'abord un court texte qui utilise naturellement presque tout le vocabulaire de cette leçon. Soulignez mentalement les mots qui parlent de participation et de temps.

\begin{textebox}[A Busy but Rewarding Term — original reading text]
Last term, I decided to get involved in several school activities. I joined the debate club and volunteered for the clean-up campaign organised by the students' council. At first, I thought I could attend every meeting, but I soon realised I had to manage my time carefully. I am currently very busy: I have classes in the morning, a part-time job at my uncle's bookshop in the afternoon, and choir practice twice a week.

To avoid wasting time, I drew up a weekly timetable and set clear priorities. Homework comes first, then community service, then sports. This organisation requires discipline and punctuality, but it is truly rewarding. Last month, our team met an important deadline: we finished the school magazine on time, and our headmaster praised our commitment and teamwork.

Of course, some activities are time-consuming, and sometimes I feel I am running out of time. But I have learned to make time for what really matters. Taking part in these activities has taught me responsibility and leadership. My advice to every student is simple: do not just watch others — participate, contribute, and you will achieve more than you ever imagined.
\end{textebox}

\textbf{Glossaire} : \emph{to get involved in} (s'impliquer dans), \emph{clean-up campaign} (campagne de nettoyage), \emph{to draw up} (établir, rédiger ; prétérit \emph{drew up}), \emph{deadline} (date limite), \emph{rewarding} (gratifiant), \emph{time-consuming} (chronophage), \emph{to run out of time} (manquer de temps).

\textbf{Questions de compréhension} : (1) Which activities did the narrator join? (2) How does he or she avoid wasting time? (3) What qualities did these activities develop?

\begin{corrige}[Questions de compréhension]
1. \eng{He/She joined the debate club and volunteered for the clean-up campaign.}
2. \eng{He/She drew up a weekly timetable and set clear priorities.}
3. \eng{They developed responsibility, leadership, discipline and teamwork.}
\end{corrige}

\courssub{Le champ lexical de la participation}

\begin{definition}[Participer : sept verbes à connaître]
Pour exprimer la participation, l'anglais dispose de plusieurs verbes proches mais pas interchangeables. \eng{To take part in} est le plus général ; \eng{to join} signifie « devenir membre de » ; \eng{to participate in} est plus soutenu ; \eng{to volunteer} insiste sur le caractère volontaire et gratuit ; \eng{to get involved in} marque l'implication personnelle ; \eng{to contribute to} signifie « apporter sa contribution à » ; \eng{to attend} signifie « être présent à, assister à ».
\end{definition}

\begin{center}
\begin{tabularx}{\linewidth}{|X|X|X|}
\hline
\rowcolor{popblueL} English & Français & Exemple \\
\hline
to take part in & participer à & She took part in the sports day. \\
\hline
to join & rejoindre, adhérer à & He joined the debate club in September. \\
\hline
to participate in & participer à (soutenu) & All students participated in the project. \\
\hline
to volunteer (for) & se porter volontaire & They volunteered for the clean-up campaign. \\
\hline
to get involved in & s'impliquer dans & More parents should get involved in school life. \\
\hline
to contribute to & contribuer à & Everyone contributed to the success of the event. \\
\hline
to attend & assister à, être présent à & I attend choir practice every Friday. \\
\hline
\end{tabularx}
\end{center}

\begin{propriete}[La préposition fait la différence]
Certains de ces verbes exigent une préposition, d'autres non ; c'est une source fréquente d'erreurs au Bac :
\begin{itemize}
\item \eng{to take part \textbf{in}}, \eng{to participate \textbf{in}}, \eng{to get involved \textbf{in}}, \eng{to contribute \textbf{to}} : la préposition est obligatoire devant le complément.
\item \eng{to join} et \eng{to attend} se construisent \textbf{sans préposition} : \eng{He joined the club.} (jamais \emph{joined to the club}) ; \eng{She attended the meeting.} (jamais \emph{attended to the meeting}).
\end{itemize}
\end{propriete}

\begin{attention}[To attend ou to assist ?]
Faux ami classique : « assister à une réunion » se dit \eng{to attend a meeting}, jamais \emph{to assist a meeting}. Le verbe \eng{to assist} existe, mais il signifie « aider » : \eng{The nurse assisted the doctor.} (« L'infirmière a aidé le médecin. »). De même, \eng{an assistant} est un « assistant/adjoint », pas un « spectateur ».
\end{attention}

\courssub{Le champ lexical des activités}

Voici les activités les plus souvent citées dans les compositions du thème \emph{Economic Life and Occupations}. Apprenez-les avec leur article et leur verbe habituel.

\begin{center}
\begin{tabularx}{\linewidth}{|X|X|}
\hline
\rowcolor{popblueL} English & Français \\
\hline
community service & le service communautaire, les travaux communautaires \\
\hline
a clean-up campaign & une campagne de nettoyage \\
\hline
sports day & la journée sportive \\
\hline
a debate club & un club de débat \\
\hline
the school choir & la chorale de l'école \\
\hline
group work & le travail de groupe \\
\hline
a part-time job & un emploi à temps partiel \\
\hline
an internship & un stage \\
\hline
\end{tabularx}
\end{center}

\eng{During my internship at a bank in Douala, I learned punctuality and discipline.} « Pendant mon stage dans une banque à Douala, j'ai appris la ponctualité et la discipline. »

\courssub{Le champ lexical de la gestion du temps}

\begin{definition}[Gérer son temps : les verbes clés]
\eng{To manage one's time} (gérer son temps), \eng{to plan} (planifier), \eng{to organise} (organiser), \eng{to set priorities} (fixer des priorités), \eng{to meet a deadline} (respecter une date limite), \eng{to avoid wasting time} (éviter de perdre du temps). Les noms correspondants sont \eng{a timetable} (emploi du temps, surtout scolaire, GB) et \eng{a schedule} (programme, planning, surtout US ; prononcé « ché-dioul » en anglais britannique).
\end{definition}

\eng{She plans her week every Sunday evening and always meets her deadlines.} « Elle planifie sa semaine chaque dimanche soir et respecte toujours ses délais. »

\begin{attention}[On time ou in time ?]
\eng{On time} signifie « à l'heure prévue, ponctuellement » : \eng{The meeting started on time.} (« La réunion a commencé à l'heure. »). \eng{In time} signifie « à temps, avant qu'il ne soit trop tard » : \eng{We arrived in time to catch the bus.} (« Nous sommes arrivés à temps pour attraper le bus. »). Ne les confondez pas dans vos compositions.
\end{attention}

\courssub{Noms et adjectifs utiles}

\textbf{Noms de qualités et de valeurs} : \eng{responsibility} (responsabilité), \eng{commitment} (engagement), \eng{teamwork} (travail d'équipe), \eng{leadership} (leadership, capacité à diriger), \eng{punctuality} (ponctualité), \eng{discipline} (discipline), \eng{achievement} (réussite, accomplissement).

\textbf{Adjectifs pour décrire les personnes et les activités} :

\begin{center}
\begin{tabularx}{\linewidth}{|X|X|X|}
\hline
\rowcolor{popblueL} Adjectif & Français & Exemple \\
\hline
busy & occupé & I am very busy this term. \\
\hline
organised & organisé & An organised student never panics. \\
\hline
punctual & ponctuel & Be punctual for every meeting. \\
\hline
efficient & efficace & She is an efficient team leader. \\
\hline
committed & engagé, dévoué & Our club members are very committed. \\
\hline
rewarding & gratifiant & Community service is a rewarding experience. \\
\hline
time-consuming & chronophage & Group work can be time-consuming. \\
\hline
\end{tabularx}
\end{center}

\courssub{Les collocations avec « time »}

\begin{definition}[Qu'est-ce qu'une collocation ?]
Une \cle{collocation} est une combinaison habituelle de mots dans une langue : les mots « vont ensemble » naturellement. On dit \eng{to make time for something} (« trouver du temps pour quelque chose ») et non \emph{to do time for something}. Les collocations ne se traduisent pas mot à mot : elles s'apprennent par cœur, comme des blocs.
\end{definition}

\begin{exemplebox}[Six collocations indispensables avec « time »]
\begin{itemize}
\item \eng{to make time for} : trouver du temps pour — \eng{I always make time for my studies.} « Je trouve toujours du temps pour mes études. »
\item \eng{to save time} : gagner du temps — \eng{Planning your day saves time.} « Planifier sa journée fait gagner du temps. »
\item \eng{to waste time} : perdre/gaspiller du temps — \eng{We must not waste time.} « Nous ne devons pas gaspiller le temps. »
\item \eng{to spend time on} : consacrer du temps à — \eng{He spends two hours on homework every day.} « Il consacre deux heures aux devoirs chaque jour. »
\item \eng{to run out of time} : manquer de temps — \eng{Hurry up! We are running out of time.} « Dépêchez-vous ! Nous manquons de temps. »
\item \eng{to keep to a schedule} : respecter un programme — \eng{If you keep to a schedule, you will finish on time.} « Si tu respectes un programme, tu finiras à temps. »
\end{itemize}
\end{exemplebox}

\begin{attention}[Actually ne veut pas dire « actuellement »]
Faux ami redoutable : \eng{I am actually busy} signifie « je suis \emph{en fait} occupé », pas « je suis actuellement occupé ». Pour « actuellement », dites \eng{currently} ou \eng{at present} : \eng{I am currently busy.} Autre piège : « une opportunité » se dit \eng{an opportunity} ; \eng{an occasion} désigne une occasion au sens d'événement (\eng{a special occasion} = une grande occasion). Enfin, \eng{She attended the community meeting.} = « Elle a assisté à la réunion communautaire » — pas « elle a aidé la réunion » !
\end{attention}

\courssub{Les familles de mots : un verbe, trois mots}

Connaître les familles de mots permet de varier les formulations dans une composition et de réussir les exercices de transformation grammaticale du Bac.

\begin{center}
\begin{tabular}{|l|l|l|}
\hline
\rowcolor{popblueL} Verbe & Nom & Adjectif / personne \\
\hline
to organise & organisation & organised \\
\hline
to participate & participation & participant \\
\hline
to manage & management & manager \\
\hline
to lead & leadership & leader \\
\hline
to achieve & achievement & achievable \\
\hline
to commit & commitment & committed \\
\hline
\end{tabular}
\end{center}

\eng{Good time management helps a manager stay organised.} « Une bonne gestion du temps aide un responsable à rester organisé. » Remarquez comment une seule famille (\emph{manage}) donne trois mots différents selon la fonction dans la phrase.

\courssub{Carte mentale : tout le vocabulaire en un coup d'œil}

\begin{popfigure}
\begin{tikzpicture}[
  mindmap,
  grow cyclic,
  every node/.style={concept, align=center, font=\small},
  concept color=popblue,
  level 1/.append style={level distance=3.4cm, sibling angle=90},
  level 2/.append style={level distance=2.4cm, font=\scriptsize}
]
\node[concept color=popblue, text=white, font=\small\bfseries] {Community \& Time}
  child[concept color=poporange] { node {Participating}
    child { node[concept color=poporangeL] {take part in, join} }
    child { node[concept color=poporangeL] {volunteer, attend} }
  }
  child[concept color=popgreen] { node {Activities}
    child { node[concept color=popgreenL] {community service, internship} }
    child { node[concept color=popgreenL] {debate club, sports day} }
  }
  child[concept color=poppurple] { node {Managing time}
    child { node[concept color=poppurpleL] {plan, set priorities} }
    child { node[concept color=poppurpleL] {meet a deadline, timetable} }
  }
  child[concept color=popgold] { node {Qualities}
    child { node[concept color=popgoldL] {punctual, committed} }
    child { node[concept color=popgoldL] {teamwork, leadership} }
  };
\end{tikzpicture}
\legende{Carte mentale du vocabulaire : quatre familles autour du thème « participation et gestion du temps ».}
\end{popfigure}

\courssub{Exercice résolu et entraînement}

\begin{exoresolu}[Choisissez le mot correct]
Complétez avec le mot qui convient : \emph{attended / actually / currently / opportunity / wasted / make}.

1. I am \ldots\ldots very busy with my internship, so I cannot join the choir this term.
2. She \ldots\ldots the community meeting last Saturday.
3. This debate competition is a great \ldots\ldots to improve your English.
4. He \ldots\ldots two hours watching videos instead of doing his homework.
5. Even when I am tired, I always \ldots\ldots time for my family.
\tcblower
\textbf{Solution détaillée} :
1. \eng{currently} — le sens est « actuellement » ; \emph{actually} signifierait « en fait » (faux ami).
2. \eng{attended} — « assister à » se dit \emph{to attend} ; \emph{to assist} signifie « aider ».
3. \eng{opportunity} — « une opportunité » ; \emph{an occasion} désigne un événement particulier.
4. \eng{wasted} — collocation : \emph{to waste time} ; on ne dit pas \emph{to lose time} pour « gaspiller ».
5. \eng{make} — collocation : \emph{to make time for} ; jamais \emph{to do time for}.
\end{exoresolu}

\begin{exercice}[À vous de jouer]
1. Traduisez en anglais : « Je participe au club de débat et je fais du bénévolat pour la campagne de nettoyage. »
2. Complétez avec un mot de la famille de \emph{organise} : « Good \ldots\ldots is the key to success. A well-\ldots\ldots student never misses a deadline. »
3. Trouvez la collocation correcte : (a) to do / make time for ; (b) to lose / waste time ; (c) to respect / meet a deadline.
4. Rédigez trois phrases sur votre emploi du temps avec \emph{timetable}, \emph{priorities} et \emph{rewarding}.
\end{exercice}

\begin{corrige}[Exercice : vocabulaire]
1. \eng{I take part in the debate club and I volunteer for the clean-up campaign.}
2. \eng{organisation} (nom, sujet de la phrase) ; \eng{organised} (adjectif devant \emph{student}).
3. (a) \eng{to make time for} ; (b) \eng{to waste time} ; (c) \eng{to meet a deadline}.
4. Modèle : \eng{My timetable is full, but I set clear priorities: school first. Managing my time well is difficult but rewarding.}
\end{corrige}

\begin{savaistu}[Le community service en Afrique anglophone]
Au Nigeria et au Ghana, le \eng{community service} fait officiellement partie du cursus scolaire : les élèves nettoient leur quartier, aident dans les bibliothèques ou plantent des arbres avant d'obtenir leur diplôme. Cela ressemble beaucoup aux travaux communautaires pratiqués dans de nombreux lycées camerounais, de Yaoundé à Bamenda. Citer ce parallèle dans votre composition montre à l'examinateur que vous reliez l'anglais au monde réel — et c'est exactement ce que la grille d'évaluation récompense.
\end{savaistu}

\begin{aretenir}[L'essentiel de la section]
\begin{itemize}
\item Participer : \eng{to take part in, to join, to volunteer, to attend} — attention, \eng{to assist} = aider.
\item Prépositions : \eng{to take part in, to participate in, to contribute to}, mais \eng{to join} et \eng{to attend} sans préposition.
\item Gérer son temps : \eng{to manage one's time, to set priorities, to meet a deadline}.
\item Collocations à mémoriser en blocs : \eng{to make time for, to waste time, to run out of time, to keep to a schedule}.
\item Faux amis : \eng{actually} = en fait (\eng{currently} = actuellement) ; \eng{an opportunity} = une opportunité ; \eng{on time} = à l'heure, \eng{in time} = à temps.
\item Familles de mots : \emph{organise / organisation / organised}, \emph{participate / participation / participant}, \emph{manage / management / manager}.
\end{itemize}
\end{aretenir}

\coursec{Modèle commenté et erreurs à éviter}

Dans la section précédente, nous avons enrichi notre vocabulaire sur la participation et la gestion du temps : \eng{to join a club}, \eng{to take part in}, \eng{to manage one's time}, \eng{a timetable}, \eng{to run out of time}... Il est maintenant temps de voir comment tous ces outils — structure, connecteurs, temps verbaux, lexique — s'assemblent dans une \cle{composition réussie}. Nous allons d'abord lire un modèle complet, le décortiquer phrase par phrase, puis passer en revue les cinq erreurs les plus fréquentes des candidats francophones, afin que vous ne les commettiez jamais.

\courssub{Lecture d'un modèle de composition}

Lisez d'abord le texte en entier, sans dictionnaire, pour en saisir l'organisation générale. Repérez mentalement les trois grandes parties : l'introduction, le développement en deux paragraphes, et la conclusion.

\begin{textebox}[Model composition — “Taking part and managing time in school activities”]
Have you ever wondered why some students do well in everything? The answer is simple: they take part in school activities and manage their time wisely.

First of all, joining clubs such as the debate club or the school choir builds confidence and teamwork. For example, last year I joined the football team, and I learned to cooperate with others.

Moreover, good time management is essential. Students who plan their week never run out of time for revision. However, some pupils waste time on their phones and then fail their exams. Therefore, making a timetable is a good solution.

In conclusion, taking part in activities and managing time are the keys to success at school and in life.
\end{textebox}

\begin{definition}[Glossaire du modèle]
\begin{itemize}
\item \eng{to wonder} (v.) : se demander
\item \eng{wisely} (adv.) : sagement, judicieusement
\item \eng{to build} (v.) : développer, construire (ici : la confiance) — irrégulier : \eng{build -- built -- built}
\item \eng{teamwork} (n.) : travail d'équipe
\item \eng{to cooperate} (v.) : coopérer, collaborer
\item \eng{to run out of time} (exp.) : manquer de temps
\item \eng{to waste time} (exp.) : perdre son temps
\item \eng{to fail an exam} (exp.) : échouer à un examen
\item \eng{the key to success} (exp.) : la clé du succès
\end{itemize}
\end{definition}

\courssub{Repérage des éléments de structure dans le modèle}

Observons comment le modèle applique exactement le plan \emph{introduction — développement — conclusion} étudié dans les sections précédentes. Chaque phrase joue un rôle précis : c'est ce rôle que le correcteur du Bac cherche dans votre copie.

\begin{center}
\begin{tabularx}{\linewidth}{|l|X|}
\hline
\rowcolor{popblueL} \textbf{Élément} & \textbf{Phrase du modèle} \\
\hline
\cle{Accroche} (question) & \eng{Have you ever wondered why some students do well in everything?} \\
\hline
\cle{Annonce de plan} & \eng{The answer is simple: they take part in school activities and manage their time wisely.} \\
\hline
\cle{Topic sentence} §1 & \eng{First of all, joining clubs... builds confidence and teamwork.} \\
\hline
\cle{Exemple} §1 & \eng{For example, last year I joined the football team...} \\
\hline
\cle{Topic sentence} §2 & \eng{Moreover, good time management is essential.} \\
\hline
\cle{Contraste} §2 & \eng{However, some pupils waste time on their phones...} \\
\hline
\cle{Conséquence / solution} & \eng{Therefore, making a timetable is a good solution.} \\
\hline
\cle{Conclusion} & \eng{In conclusion, taking part in activities and managing time are the keys to success...} \\
\hline
\end{tabularx}
\end{center}

\begin{aretenir}[Les trois réflexes du bon rédacteur]
\begin{enumerate}
\item Chaque paragraphe commence par une \cle{topic sentence} qui annonce l'idée principale.
\item Chaque idée est \cle{illustrée} par un exemple (\eng{For example}, \eng{For instance}).
\item Les paragraphes sont \cle{reliés} par des connecteurs variés : \eng{First of all}, \eng{Moreover}, \eng{However}, \eng{Therefore}, \eng{In conclusion}.
\end{enumerate}
\end{aretenir}

\courssub{Analyse des temps verbaux du modèle}

Le choix des temps n'est jamais le fruit du hasard : chaque temps répond à une intention. Voici la logique qui sous-tend le modèle.

\begin{propriete}[Quels temps dans une composition d'argumentation ?]
\begin{itemize}
\item Le \cle{present simple} domine : il exprime des vérités générales et des habitudes. \eng{Students who plan their week never run out of time.} — Les élèves qui planifient leur semaine ne manquent jamais de temps.
\item Le \cle{past simple} sert à raconter une expérience personnelle datée : \eng{last year I joined the football team, and I learned to cooperate.} — l'année dernière, je me suis inscrit à l'équipe de football et j'ai appris à coopérer.
\item Le \cle{present perfect} convient pour une expérience sans date précise : \eng{Have you ever wondered...?} — Vous êtes-vous déjà demandé... ?
\end{itemize}
\end{propriete}

\begin{exemplebox}[Exemples traduits tirés du modèle]
\eng{Students who plan their week never run out of time for revision.} — Les élèves qui planifient leur semaine ne manquent jamais de temps pour réviser.

\eng{Making a timetable is a good solution.} — Faire un emploi du temps est une bonne solution.

\eng{Joining clubs builds confidence and teamwork.} — Rejoindre des clubs développe la confiance et l'esprit d'équipe.
\end{exemplebox}

Remarquez la construction \eng{making a timetable is...} : en anglais, le sujet d'une phrase peut être un \cle{gérondif} (verbe en \emph{-ing}), alors que le français utilise l'infinitif (« faire un emploi du temps \emph{est}... »). Le gérondif sujet se comporte comme un singulier : \eng{Joining clubs \textbf{builds} confidence} (et non \eng{build}).

\courssub{Les cinq erreurs à éviter}

\begin{attention}[Erreur n°1 : le hors-sujet]
Rédiger deux pages sur vos loisirs (la danse, la musique, les jeux vidéo) sans jamais parler de \cle{participation} ni de \cle{gestion du temps}, c'est du hors-sujet. Chaque phrase doit répondre à la question : \emph{en quoi cette idée concerne-t-elle la participation aux activités ou la gestion du temps ?} Avant d'écrire, soulignez les mots-clés du sujet et vérifiez à la relecture que chacun apparaît dans votre texte.
\end{attention}

\begin{attention}[Erreur n°2 : les phrases interminables sans connecteurs]
Une phrase = une idée. Évitez : \eng{I like football and I play every day and my friends also play and we win matches and it is good.} Préférez : \eng{I like football. I play every day with my friends. Moreover, playing in a team teaches us cooperation.} Des phrases courtes, reliées par des connecteurs, donnent un texte clair et noté plus haut.
\end{attention}

\begin{attention}[Erreur n°3 : la traduction mot à mot du français]
Ne dites jamais \eng{I have 17 years} (calque de « j'ai 17 ans ») : dites \eng{I am 17 years old}. De même : \eng{to participate \textbf{to} an activity} est faux ; dites \eng{to participate \textbf{in} an activity} ou \eng{to take part \textbf{in} an activity}. Autres calques fréquents : \eng{I am agree} → \eng{I agree} ; \eng{to make a fault} → \eng{to make a mistake}.
\end{attention}

\begin{attention}[Erreur n°4 : oublier le -s de la troisième personne]
Au present simple, la troisième personne du singulier prend un \cle{-s} : \eng{he manage\textbf{s}}, \eng{she take\textbf{s} part}, \eng{it build\textbf{s} confidence}. Attention au piège inverse : au pluriel, il n'y a \emph{pas} de -s (\eng{students waste}, et non \eng{students wastes}). C'est l'erreur la plus sanctionnée au Bac. Astuce : à la relecture, soulignez chaque sujet à la troisième personne du singulier et vérifiez le verbe.
\end{attention}

\begin{attention}[Erreur n°5 : négliger la relecture]
Une composition pleine de fautes d'accord, d'orthographe, de ponctuation ou de majuscules perd des points même si les idées sont bonnes. Réservez toujours les \cle{cinq dernières minutes} à la relecture.
\end{attention}

\begin{methode}[La relecture en trois passages]
\begin{enumerate}
\item \textbf{Premier passage — le sens et le plan} : mon introduction contient-elle une accroche et une annonce de plan ? Chaque paragraphe a-t-il une topic sentence ? Ma conclusion répond-elle au sujet ?
\item \textbf{Deuxième passage — la grammaire et les temps} : -s à la troisième personne, past simple pour les expériences datées, prépositions correctes (\eng{take part in}, \eng{interested in}, \eng{good at}).
\item \textbf{Troisième passage — l'orthographe et la ponctuation} : majuscules en début de phrase et pour \eng{I}, point final, virgule après les connecteurs en tête de phrase (\eng{Moreover, ...}), orthographe des mots difficiles (\eng{timetable}, \eng{confidence}, \eng{success}).
\end{enumerate}
\end{methode}

\begin{popfigure}
\begin{center}
\begin{tikzpicture}[font=\small]
\node[draw=popblue, fill=popblueL, rounded corners, thick, minimum width=4.2cm, minimum height=0.9cm, align=center] (t) at (0,0) {\textbf{Checklist de relecture}};
\node[draw=popgreen, fill=popgreenL, rounded corners, minimum width=4.6cm, minimum height=0.8cm, align=center] (a) at (0,-1.3) {$\square$ Verbes : temps et -s corrects};
\node[draw=poporange, fill=poporangeL, rounded corners, minimum width=4.6cm, minimum height=0.8cm, align=center] (b) at (0,-2.4) {$\square$ Connecteurs variés};
\node[draw=poppurple, fill=poppurpleL, rounded corners, minimum width=4.6cm, minimum height=0.8cm, align=center] (c) at (0,-3.5) {$\square$ Plan respecté (intro / développement / conclusion)};
\node[draw=poppink, fill=poppinkL, rounded corners, minimum width=4.6cm, minimum height=0.8cm, align=center] (d) at (0,-4.6) {$\square$ Orthographe et ponctuation};
\node[draw=popgold, fill=popgoldL, rounded corners, minimum width=4.6cm, minimum height=0.8cm, align=center] (e) at (0,-5.7) {$\square$ Nombre de mots respecté};
\draw[-{Stealth}, thick, popblue] (t) -- (a);
\draw[-{Stealth}, thick, popblue] (a) -- (b);
\draw[-{Stealth}, thick, popblue] (b) -- (c);
\draw[-{Stealth}, thick, popblue] (c) -- (d);
\draw[-{Stealth}, thick, popblue] (d) -- (e);
\end{tikzpicture}
\end{center}
\legende{La checklist de relecture : cinq cases à cocher avant de rendre sa copie.}
\end{popfigure}

\begin{center}
\begin{tabularx}{\linewidth}{|X|X|X|}
\hline
\rowcolor{popblueL} \textbf{Français (calque à éviter)} & \textbf{English (correct)} & \textbf{Remarque} \\
\hline
J'ai 17 ans. & I am 17 years old. & \eng{to be}, pas \eng{to have} \\
\hline
participer à une activité & to take part in an activity & préposition \eng{in} \\
\hline
Je suis d'accord. & I agree. & pas de \eng{am} \\
\hline
Il gère bien son temps. & He manages his time well. & -s à la 3\textsuperscript{e} personne \\
\hline
perdre son temps & to waste time & pas \eng{to lose time} \\
\hline
réussir à un examen & to pass an exam & attention : « passer un examen » se dit \eng{to sit (for) an exam} ou \eng{to take an exam} \\
\hline
\end{tabularx}
\end{center}

\begin{exoresolu}[Corriger un extrait de copie d'élève]
Énoncé — Voici un extrait de la copie de Marlyse, élève de Tle A à Buea. Relevez et corrigez les erreurs : \eng{I have 17 years and I participate to the debate club. Last year I join the club and I learn many things. My friend she manage her time very well, so she never fail her exams.}
\tcblower
Solution détaillée :
\begin{enumerate}
\item \eng{I have 17 years} → \eng{I am 17 years old} (calque du français « j'ai 17 ans »).
\item \eng{I participate to} → \eng{I participate \textbf{in}} ou \eng{I take part in} (préposition \eng{in}).
\item \eng{Last year I join} → \eng{Last year I \textbf{joined}} (expérience datée : past simple).
\item \eng{I learn many things} → \eng{I \textbf{learned} many things} (même raison : le repère \eng{last year} impose le passé).
\item \eng{My friend she manage} → \eng{My friend \textbf{manages}} (double sujet à supprimer + -s de la troisième personne).
\item \eng{she never fail} → \eng{she never \textbf{fails}} (-s de la troisième personne).
\end{enumerate}
Version corrigée : \eng{I am 17 years old and I take part in the debate club. Last year I joined the club and I learned many things. My friend manages her time very well, so she never fails her exams.}
\end{exoresolu}

\begin{exercice}[À vous de jouer]
Réécrivez le paragraphe suivant en corrigeant les erreurs et en ajoutant deux connecteurs : \eng{I am agree that time management is important. Students wastes time on phones. Last term I make a timetable and my results improve. Participate to school activities is also good.}
\end{exercice}

\begin{corrige}[Exercice — Réécriture]
\eng{I agree that time management is important. However, many students waste time on their phones. Last term I made a timetable and my results improved. Moreover, taking part in school activities is also beneficial.}
Corrections : \eng{I am agree} → \eng{I agree} ; \eng{students wastes} → \eng{students waste} (sujet pluriel : pas de -s) ; \eng{I make} → \eng{I made} (past simple, verbe irrégulier \eng{make -- made -- made}) ; \eng{improve} → \eng{improved} ; \eng{participate to} → \eng{taking part in} (gérondif sujet + préposition \eng{in}) ; ajout des connecteurs \eng{However} et \eng{Moreover}.
\end{corrige}

\begin{aretenir}[L'essentiel de la section]
Un bon modèle se reconnaît à trois signes : un plan visible (accroche, annonce, topic sentences, conclusion), des connecteurs variés et des temps verbaux justifiés (present simple pour les généralités, past simple pour les expériences datées, present perfect pour les expériences non datées). Les cinq erreurs à bannir : le hors-sujet, les phrases trop longues, les calques du français, l'oubli du -s (ou le -s abusif au pluriel), et l'absence de relecture. Gardez la checklist en tête : verbes, connecteurs, plan, orthographe, nombre de mots.
\end{aretenir}

\coursec{Grille d'évaluation et entraînement au Bac}

Dans la section précédente, nous avons étudié un modèle commenté de composition et repéré les erreurs les plus fréquentes des candidats francophones. Il reste une dernière étape, décisive : apprendre à \cle{juger} une copie comme le ferait un correcteur du Baccalauréat, puis à s'entraîner dans les conditions réelles de l'examen. C'est l'objet de cette dernière section : connaître la grille d'évaluation, s'auto-évaluer, évaluer ses pairs et rédiger en 45 minutes.

\courssub{Pourquoi évaluer avant de rendre sa copie ?}

Observez d'abord ces deux extraits de copies corrigées sur le sujet \eng{“Describe how you take part in school activities and manage your time.”}

\begin{exemplebox}[Deux copies, deux destins]
\textbf{Copy A :} \eng{“I like school. I do many activities. I play football. I am in the choir. I have no time but I try.”} — Pas d'introduction, pas de connecteurs, phrases courtes et répétitives. Note obtenue : 8/20.

\textbf{Copy B :} \eng{“First of all, I take part in several school activities, such as the football team and the English club. Moreover, I manage my time carefully: every evening, I plan the next day in a small notebook. In conclusion, taking part in activities has taught me discipline and punctuality.”} — Structure visible, connecteurs variés, vocabulaire riche. Note obtenue : 16/20.
\end{exemplebox}

La différence entre les deux copies n'est pas seulement une question de niveau d'anglais : c'est surtout une question de \cle{stratégie}. La copie B répond aux cinq critères que tout correcteur applique, consciemment ou non. Apprendre ces critères, c'est apprendre à écrire \emph{pour être bien noté}.

\courssub{La grille d'évaluation type}

\begin{propriete}[Grille d'évaluation d'une composition (sur 20)]
\begin{center}
\begin{tabularx}{\linewidth}{|X|c|}
\hline
\rowcolor{popblueL} \textbf{Critère} & \textbf{Points} \\
\hline
\cle{Pertinence} du contenu : le sujet est respecté, la consigne est suivie, les idées répondent à la question posée. & /4 \\
\hline
\cle{Organisation} et connecteurs : introduction, paragraphes distincts, conclusion ; connecteurs logiques visibles et variés. & /4 \\
\hline
\cle{Vocabulaire} : lexique riche, précis, adapté au thème (participation, gestion du temps) ; pas de répétitions. & /4 \\
\hline
\cle{Grammaire} et temps : temps verbaux corrects et cohérents, accords, 3e personne du singulier (\eng{he takes}, \eng{she manages}). & /4 \\
\hline
\cle{Orthographe} et présentation : spelling, ponctuation, majuscules, alinéas, écriture lisible. & /4 \\
\hline
\textbf{Total} & \textbf{/20} \\
\hline
\end{tabularx}
\end{center}
\end{propriete}

Détaillons chaque critère.

\textbf{1. Pertinence (respect du sujet).} Le correcteur se demande : \eng{“Did the candidate answer the question?”} Si le sujet demande de parler de la gestion du temps \emph{et} de la participation aux activités, une copie qui ne traite que du football perd des points, même si l'anglais est excellent. Soulignez les mots-clés du sujet avant d'écrire et vérifiez à la fin que chacun a été traité.

\textbf{2. Organisation.} Une composition doit montrer \emph{visuellement} sa structure : un alinéa pour l'introduction, un alinéa par grande idée, un alinéa pour la conclusion. Les connecteurs (\eng{First of all}, \eng{Moreover}, \eng{However}, \eng{In conclusion}) sont les panneaux indicateurs de votre texte.

\textbf{3. Vocabulaire.} Évitez les répétitions : au lieu d'écrire trois fois \eng{activities}, variez avec \eng{events}, \eng{clubs}, \eng{duties}, \eng{tasks}. Montrez que vous connaissez le lexique du thème : \eng{to take part in}, \eng{to volunteer}, \eng{a deadline}, \eng{a schedule}, \eng{punctual}.

\textbf{4. Grammaire.} Les trois fautes les plus sanctionnées : l'oubli du \emph{-s} à la 3e personne du singulier (\eng{he take} au lieu de \eng{he takes}), le mélange incohérent des temps (passé et présent mélangés sans raison), et les accords (\eng{this activities} au lieu de \eng{these activities}).

\textbf{5. Orthographe et présentation.} Une copie lisible, avec des alinéas nets, des majuscules en début de phrase et une ponctuation correcte, donne immédiatement une bonne impression. Les mots mal orthographiés (\eng{wich} pour \eng{which}, \eng{recieve} pour \eng{receive}) coûtent des points.

\begin{attention}[Pièges fréquents des candidats francophones]
\begin{itemize}
\item \eng{I am agree} : faux ami avec « je suis d'accord ». En anglais : \eng{I agree} (verbe simple, sans \emph{be}).
\item \eng{Actually} ne signifie pas « actuellement » mais « en fait ». « Actuellement » se dit \eng{currently} ou \eng{at present}.
\item \eng{An advice} : \emph{advice} est indénombrable. Dites \eng{some advice} ou \eng{a piece of advice}.
\item \eng{I have 17 years} : traduction mot à mot de « j'ai 17 ans ». En anglais : \eng{I am 17 (years old)}.
\item Ne mélangez pas les temps : si vous racontez une expérience passée, restez au prétérit ; si vous décrivez vos habitudes, restez au présent simple.
\end{itemize}
\end{attention}

\courssub{Figure : la grille d'évaluation en un coup d'œil}

\begin{popfigure}
\begin{tikzpicture}[font=\small]
\draw[fill=popblue, draw=popblue] (0,0) rectangle (10.5,0.8);
\node[text=white, font=\bfseries] at (5.25,0.4) {Evaluation grid -- Grille d'évaluation};
\draw[fill=popblueL, draw=popline] (0,-0.8) rectangle (10.5,0);
\node[anchor=west, font=\bfseries] at (0.2,-0.4) {Criterion};
\node[font=\bfseries] at (7.3,-0.4) {Points};
\node[anchor=west, font=\bfseries] at (8.1,-0.4) {Comments};
\draw[fill=white, draw=popline] (0,-1.6) rectangle (10.5,-0.8);
\node[anchor=west] at (0.2,-1.2) {1. Relevance (Pertinence)};
\node at (7.3,-1.2) {/4};
\node[anchor=west, text=popmuted] at (8.1,-1.2) {topic respected?};
\draw[fill=popcream, draw=popline] (0,-2.4) rectangle (10.5,-1.6);
\node[anchor=west] at (0.2,-2.0) {2. Organisation \& connectors};
\node at (7.3,-2.0) {/4};
\node[anchor=west, text=popmuted] at (8.1,-2.0) {intro / body / conclusion?};
\draw[fill=white, draw=popline] (0,-3.2) rectangle (10.5,-2.4);
\node[anchor=west] at (0.2,-2.8) {3. Vocabulary};
\node at (7.3,-2.8) {/4};
\node[anchor=west, text=popmuted] at (8.1,-2.8) {rich, varied, no repetition?};
\draw[fill=popcream, draw=popline] (0,-4.0) rectangle (10.5,-3.2);
\node[anchor=west] at (0.2,-3.6) {4. Grammar \& tenses};
\node at (7.3,-3.6) {/4};
\node[anchor=west, text=popmuted] at (8.1,-3.6) {3rd person -s, tense agreement?};
\draw[fill=white, draw=popline] (0,-4.8) rectangle (10.5,-4.0);
\node[anchor=west] at (0.2,-4.4) {5. Spelling \& presentation};
\node at (7.3,-4.4) {/4};
\node[anchor=west, text=popmuted] at (8.1,-4.4) {paragraphs, punctuation, legibility?};
\draw[fill=poporangeL, draw=popline] (0,-5.6) rectangle (10.5,-4.8);
\node[anchor=west, font=\bfseries] at (0.2,-5.2) {TOTAL};
\node[font=\bfseries] at (7.3,-5.2) {/20};
\end{tikzpicture}
\legende{La grille d'évaluation type : cinq critères, quatre points chacun, pour un total de vingt points.}
\end{popfigure}

\courssub{Auto-évaluation : le code COCO}

Avant de rendre votre copie, relisez-la en suivant un code mnémotechnique simple : \cle{COCO}.

\begin{methode}[S'auto-corriger avec le code COCO]
\begin{enumerate}
\item \textbf{C -- Content (contenu).} Ai-je répondu à la question posée ? Ai-je traité tous les mots-clés du sujet ? Chaque paragraphe apporte-t-il une idée en lien avec le thème ?
\item \textbf{O -- Organisation.} Mon introduction est-elle visible ? Y a-t-il un alinéa par idée ? Ma conclusion est-elle clairement marquée (\eng{In conclusion}, \eng{To sum up}) ? Ai-je utilisé au moins quatre connecteurs différents ?
\item \textbf{C -- Correction (grammaire).} Ai-je vérifié le \emph{-s} de la 3e personne du singulier ? Mes temps sont-ils cohérents ? Mes accords (this/these, a/an) sont-ils corrects ?
\item \textbf{O -- Orthography (orthographe).} Ai-je relu mot par mot pour traquer les fautes de spelling ? Majuscules en début de phrase ? Ponctuation ? Écriture lisible et alinéas nets ?
\end{enumerate}
Réservez toujours les cinq dernières minutes de l'épreuve à cette relecture COCO : c'est là que se gagnent deux à trois points.
\end{methode}

\courssub{Évaluation par les pairs}

L'évaluation par les pairs (\eng{peer assessment}) consiste à échanger les copies et à annoter celle de son voisin avec la grille. La règle d'or : formuler \cle{deux points forts} et \cle{un conseil} (\eng{two strengths and one piece of advice}). Le commentaire doit être précis, bienveillant et utile.

\begin{exemplebox}[Commentaires de pairs : modèles à imiter]
\eng{“Your introduction is clear, but add more connectors in paragraph 2.”} — « Ton introduction est claire, mais ajoute plus de connecteurs au paragraphe 2. »

\eng{“Good vocabulary about time management; however, check the third person -s: ‘he take' should be ‘he takes'.”} — « Bon vocabulaire sur la gestion du temps ; cependant, vérifie le -s de la 3e personne : “he take” devrait être “he takes”. »

\eng{“I liked your conclusion. Be careful with spelling: ‘wich' is wrong -- write ‘which'.”} — « J'ai aimé ta conclusion. Attention à l'orthographe : “wich” est faux -- écris “which”. »

\eng{“Your ideas are relevant, but write shorter paragraphs and start a new line for each idea.”} -- « Tes idées sont pertinentes, mais fais des paragraphes plus courts et va à la ligne pour chaque idée. »
\end{exemplebox}

\begin{center}
\begin{tabularx}{\linewidth}{|X|X|X|}
\hline
\rowcolor{popblueL} \textbf{Français} & \textbf{English} & \textbf{Remarque} \\
\hline
Ton introduction est claire. & Your introduction is clear. & point fort \\
\hline
Ajoute plus de connecteurs. & Add more connectors. & conseil (impératif) \\
\hline
Vérifie les temps. & Check your tenses. & conseil grammatical \\
\hline
Attention à l'orthographe. & Be careful with spelling. & \emph{spelling} indénombrable \\
\hline
Ta conclusion est convaincante. & Your conclusion is convincing. & point fort \\
\hline
Fais des paragraphes plus courts. & Write shorter paragraphs. & conseil de présentation \\
\hline
\end{tabularx}
\end{center}

\begin{experience}[Atelier d'évaluation par les pairs]
\begin{enumerate}
\item Rédigez individuellement une composition de 150 à 180 mots sur le sujet : \eng{“How do you take part in community activities and manage your time?”}
\item Échangez votre copie avec votre voisin.
\item Annotez la copie reçue avec la grille : attribuez une note sur 4 pour chacun des cinq critères, puis écrivez en marge deux points forts et un conseil, en anglais.
\item Rendez la copie à son auteur et discutez ensemble des annotations.
\item Réécrivez votre propre composition en tenant compte du conseil reçu.
\end{enumerate}
\end{experience}

\courssub{Entraînement chronométré : les conditions du Bac}

Le jour de l'examen, vous disposerez d'environ 45 minutes pour votre composition. Voici comment répartir ce temps.

\begin{methode}[Gérer ses 45 minutes à l'examen]
\begin{enumerate}
\item \textbf{5 minutes -- Analyse du sujet.} Soulignez les mots-clés, identifiez le type de texte attendu et le temps verbal dominant.
\item \textbf{5 minutes -- Plan.} Notez au brouillon : une idée pour l'introduction, deux ou trois idées pour le développement, une idée pour la conclusion, plus quatre ou cinq connecteurs à placer.
\item \textbf{25 minutes -- Rédaction.} Écrivez proprement, un alinéa par idée, sans ratures excessives.
\item \textbf{10 minutes -- Relecture COCO.} Content, Organisation, Correction, Orthography. Corrigez le \emph{-s} de la 3e personne, les temps, l'orthographe.
\end{enumerate}
\end{methode}

\begin{savaistu}[Le secret des copies bien notées]
Au Baccalauréat, chaque correcteur lit des centaines de copies en quelques jours. Une composition bien paragraphe, avec des connecteurs immédiatement visibles (\eng{First of all}, \eng{Moreover}, \eng{However}, \eng{In conclusion}), est repérée en quelques secondes comme une copie sérieuse -- et elle est presque toujours mieux notée qu'une copie au contenu équivalent mais présentée en un seul bloc. La forme n'est pas un détail : elle fait partie de la compétence évaluée. Écrire lisiblement, sauter une ligne entre les paragraphes et soigner la première phrase de l'introduction sont des gestes qui rapportent des points.
\end{savaistu}

\begin{exoresolu}[Évaluer une copie avec la grille]
\textbf{Énoncé.} Voici la composition de Manka, élève de Tle A à Buea, sur le sujet \eng{“Describe how you take part in school activities and manage your time.”} Notez-la sur 20 avec la grille et rédigez deux points forts et un conseil.

\eng{“I take part in many activity at school. I am member of the football team and the debate club. Every day I wake up at 5 o'clock and I plan my day. First I do my homework, after I go to training. Last week, our team win a match against a school from Douala. I am very happy because I manage my time good.”}

\tcblower
\textbf{Solution détaillée.}

\textbf{Pertinence : 3/4.} Le sujet est globalement respecté (participation + gestion du temps), mais le développement est mince : une seule phrase sur la gestion du temps.

\textbf{Organisation : 1/4.} Aucun alinéa, pas d'introduction ni de conclusion identifiables ; seuls deux connecteurs (\eng{First}, \eng{after}).

\textbf{Vocabulaire : 3/4.} Lexique correct (\eng{football team}, \eng{debate club}, \eng{training}), mais répétitions et pauvreté relative.

\textbf{Grammaire : 1/4.} Fautes : \eng{many activity} (au lieu de \eng{many activities}), \eng{I am member} (il manque \emph{a}), \eng{our team win} (prétérit : \eng{won}), \eng{good} (adverbe : \eng{well}).

\textbf{Orthographe et présentation : 2/4.} Orthographe correcte mais aucun alinéa, texte en un seul bloc.

\textbf{Total : 10/20.}

\textbf{Deux points forts :} \eng{“Your ideas are relevant and you give a concrete example (the match against Douala). Your spelling is good.”}

\textbf{Un conseil :} \eng{“Organise your text into clear paragraphs with an introduction and a conclusion, and check your tenses and agreements: ‘many activities', ‘a member', ‘won', ‘well'.”}
\end{exoresolu}

\begin{exercice}[Entraînement chronométré -- conditions d'examen]
En 45 minutes, sans dictionnaire, rédigez une composition de 180 à 200 mots sur le sujet suivant :

\eng{“Your school is organising a Community Day. Write a composition explaining how you will take part in the activities and how you will manage your time before, during and after the event.”}

Avant de rendre votre copie, appliquez la relecture COCO et cochez chaque critère de la grille. Échangez ensuite votre copie avec un camarade pour une évaluation croisée (deux points forts, un conseil).
\end{exercice}

\begin{corrige}[Exercice -- éléments de correction]
Il n'existe pas de copie unique, mais toute bonne réponse doit contenir :
\begin{itemize}
\item Une \textbf{introduction} qui présente le Community Day et annonce le plan : \eng{“Our school is organising a Community Day next month. In this composition, I will explain how I intend to take part and how I will manage my time.”}
\item Un \textbf{développement} en deux ou trois paragraphes : les activités choisies (\eng{cleaning the school yard}, \eng{visiting the orphanage}, \eng{playing in the football match}), puis l'organisation du temps (\eng{I will finish my homework the day before}, \eng{I will wake up early}, \eng{I will follow the programme}).
\item Des \textbf{connecteurs variés} : \eng{First of all}, \eng{Moreover}, \eng{In addition}, \eng{However}, \eng{Finally}, \eng{In conclusion}.
\item Des \textbf{temps cohérents} : futur (\eng{will} / \eng{going to}) pour les projets, présent simple pour les habitudes.
\item Une \textbf{conclusion} : \eng{“In conclusion, the Community Day will teach me solidarity and time management.”}
\end{itemize}
Notation : appliquez la grille -- Pertinence /4, Organisation /4, Vocabulaire /4, Grammaire /4, Orthographe et présentation /4. Une copie complète, bien paragraphe et relue avec le code COCO vise 16/20 et plus.
\end{corrige}

\begin{aretenir}[L'essentiel de la section]
\begin{itemize}
\item Une composition est notée sur cinq critères de quatre points : pertinence, organisation, vocabulaire, grammaire, orthographe et présentation.
\item Avant de rendre sa copie, on applique le code \cle{COCO} : Content, Organisation, Correction, Orthography.
\item L'évaluation par les pairs suit la règle : deux points forts + un conseil, formulés en anglais précis et bienveillant.
\item À l'examen, on répartit ses 45 minutes : 5 min d'analyse, 5 min de plan, 25 min de rédaction, 10 min de relecture.
\item La présentation (alinéas, connecteurs visibles, écriture lisible) rapporte des points : elle fait partie de la compétence évaluée.
\end{itemize}
\end{aretenir}

\newpage

\coursec{Activité d'intégration}

\courssub{Objectifs de l'activité}

À l'issue de cette activité d'intégration, l'élève sera capable de :
\begin{itemize}
\item analyser un sujet de composition et dégager un plan en trois parties (\cle{introduction}, \cle{body}, \cle{conclusion}) ;
\item rédiger, en temps limité (45 minutes), une composition de 180 à 220 mots sur la participation à une activité communautaire et la gestion du temps ;
\item employer correctement au moins \textbf{six connecteurs logiques différents} et \textbf{cinq mots du vocabulaire} de la leçon ;
\item utiliser les temps verbaux appropriés (past simple pour le récit, present simple pour les habitudes et les vérités générales, futur pour les projets) ;
\item évaluer la copie d'un camarade à l'aide d'une grille d'évaluation et formuler un retour constructif (\emph{two strengths and one piece of advice}) ;
\item relire et corriger sa propre production avant de la remettre.
\end{itemize}

\courssub{Situation}

Vous êtes membre du \emph{English Club} du lycée de Biyem-Assi, à Yaoundé. Le club organise, samedi prochain, un \emph{Community Work Day} : le matin, nettoyage du marché de Mokolo avec les vendeuses ; l'après-midi, rangement de la bibliothèque municipale et séance de lecture avec les enfants du quartier. Le proviseur a demandé au club de produire un texte en anglais qui sera publié dans le journal du lycée et affiché au tableau d'affichage du club. Voici l'annonce officielle de l'événement :

\begin{textebox}[English Club — Official Announcement]
\textbf{COMMUNITY WORK DAY — Saturday, 14th October}

Dear members,

Our English Club is organising a Community Work Day. In the morning, from 7 a.m. to 11 a.m., we shall clean Mokolo market together with the traders. In the afternoon, from 2 p.m. to 5 p.m., we shall arrange books at the municipal library and read stories with the children of the neighbourhood.

Every member must take part and manage his or her time carefully. Bring gloves, a broom and a notebook. Do not be late: punctuality is the key to success!

The best composition about this day will be published in the school magazine.

The President of the English Club
\end{textebox}

\begin{definition}[Mots difficiles de l'annonce]
\begin{itemize}
\item \eng{traders} (n.) : commerçants, vendeurs du marché.
\item \eng{to arrange books} (v.) : ranger des livres.
\item \eng{neighbourhood} (n.) : quartier. Orthographe britannique ; en anglais américain : \eng{neighborhood}.
\item \eng{punctuality} (n.) : ponctualité. Se prononce « ponk-tchou-a-li-ti ».
\item \eng{the key to success} (expr.) : la clé du succès. Attention à la préposition : \eng{the key \textbf{to} success}.
\end{itemize}
\end{definition}

Votre tâche : rédiger la composition qui sera peut-être publiée dans le journal du lycée.

\courssub{Consignes}

\textbf{Tâche 1 — Compréhension et planification (10 minutes).} Lisez attentivement l'annonce du club. Soulignez les informations essentielles : \emph{who? what? when? where? why?} Puis complétez le plan en trois parties ci-dessous avec vos idées personnelles (au brouillon) :

\begin{enumerate}
\item \textbf{Introduction} : présentez l'événement et annoncez votre sujet (\eng{Last Saturday, our English Club organised...}).
\item \textbf{Développement} : deux ou trois paragraphes — (a) ce que vous avez fait et comment vous avez participé ; (b) comment vous avez géré votre temps (emploi du temps, ponctualité, priorités) ; (c) ce que cette expérience vous a appris.
\item \textbf{Conclusion} : bilan personnel et ouverture (\eng{In conclusion, taking part in... taught me that...}).
\end{enumerate}

\textbf{Tâche 2 — Mobilisation du lexique (5 minutes).} Sur votre brouillon, choisissez au moins \textbf{cinq mots ou expressions} de la liste suivante et entourez-les : \eng{to take part in}, \eng{to volunteer}, \eng{community work}, \eng{to manage one's time}, \eng{timetable}, \eng{punctual}, \eng{to waste time}, \eng{priority}, \eng{teamwork}, \eng{to schedule}, \eng{deadline}, \eng{commitment}.

\textbf{Tâche 3 — Rédaction (45 minutes).} Rédigez votre composition intitulée \emph{“Taking part in our Community Work Day and managing my time”}, en \textbf{180 à 220 mots}. Respectez les contraintes suivantes :
\begin{enumerate}
\item trois parties nettement séparées (introduction, développement, conclusion) ;
\item au moins \textbf{six connecteurs différents} parmi : \eng{first of all}, \eng{then}, \eng{moreover}, \eng{however}, \eng{as a result}, \eng{in addition}, \eng{finally}, \eng{in conclusion}, \eng{although}, \eng{because} ;
\item au moins \textbf{cinq mots du vocabulaire} de la Tâche 2 ;
\item des temps verbaux cohérents : \cle{past simple} pour raconter la journée, \cle{present simple} pour les vérités générales, \cle{future} ou \eng{going to} pour les projets.
\end{enumerate}

\textbf{Tâche 4 — Évaluation croisée (15 minutes).} Échangez votre copie avec votre voisin. Corrigez-la à l'aide de la grille d'évaluation ci-dessous, puis rédigez votre retour en anglais : \emph{two strengths and one piece of advice} (deux points forts et un conseil).

\begin{center}
\begin{tabularx}{\linewidth}{|X|c|c|}
\hline
\rowcolor{popblueL} \textbf{Critère} & \textbf{Points} & \textbf{Obtenu} \\
\hline
Respect du sujet et de la longueur (180--220 mots) & 4 & \\
\hline
Structure claire : introduction, développement, conclusion & 4 & \\
\hline
Connecteurs logiques (au moins 6, variés et corrects) & 4 & \\
\hline
Vocabulaire de la leçon (au moins 5 mots, bien employés) & 4 & \\
\hline
Correction grammaticale (temps, accords, orthographe) & 4 & \\
\hline
\textbf{Total} & \textbf{20} & \\
\hline
\end{tabularx}
\end{center}

\textbf{Tâche 5 — Relecture finale et remise (10 minutes).} Reprenez votre copie, intégrez le conseil de votre camarade, corrigez les erreurs signalées et remettez la version finale au professeur. Les trois meilleures compositions seront lues à voix haute et affichées au tableau du club.

\begin{methode}[Relire efficacement sa composition en 5 minutes]
\begin{enumerate}
\item \textbf{Contenu} : ai-je répondu à toutes les parties du sujet ?
\item \textbf{Structure} : mes trois parties sont-elles visibles et équilibrées ?
\item \textbf{Connecteurs} : en ai-je au moins six, différents les uns des autres ?
\item \textbf{Grammaire} : vérifiez chaque verbe (temps, troisième personne du singulier, verbes irréguliers au passé).
\item \textbf{Orthographe et ponctuation} : majuscules en début de phrase, \eng{I} toujours en majuscule, pas de mot français déguisé.
\end{enumerate}
\end{methode}

\courssub{Ressources à mobiliser}

\begin{aretenir}[Connecteurs logiques par fonction]
\begin{itemize}
\item \textbf{Ordre} : \eng{first of all}, \eng{then}, \eng{after that}, \eng{finally}.
\item \textbf{Ajout} : \eng{moreover}, \eng{in addition}, \eng{furthermore}.
\item \textbf{Opposition} : \eng{however}, \eng{although}, \eng{nevertheless}.
\item \textbf{Cause et conséquence} : \eng{because}, \eng{as a result}, \eng{therefore}.
\item \textbf{Conclusion} : \eng{in conclusion}, \eng{to sum up}, \eng{all in all}.
\end{itemize}
\end{aretenir}

\begin{propriete}[Temps verbaux dans ce type de composition]
\begin{center}
\begin{tabular}{|l|l|l|}
\hline
\rowcolor{popblueL} \textbf{Temps} & \textbf{Emploi} & \textbf{Exemple} \\
\hline
Past simple & raconter la journée & We \textbf{cleaned} the market. \\
\hline
Present simple & habitudes, vérités & Punctuality \textbf{is} important. \\
\hline
Be going to / will & projets futurs & I \textbf{am going to} volunteer again. \\
\hline
\end{tabular}
\end{center}
Rappel : au past simple, les verbes irréguliers fréquents de ce thème sont \eng{meet -- met -- met}, \eng{go -- went -- gone}, \eng{take -- took -- taken}, \eng{teach -- taught -- taught}, \eng{write -- wrote -- written}.
\end{propriete}

\begin{attention}[Erreurs fréquentes des francophones]
\begin{itemize}
\item Ne dites pas \eng{*I assisted to the event} : dites \eng{I \textbf{attended} the event} ou \eng{I \textbf{took part in} the event}. \emph{Assister à} ne se traduit pas par \emph{to assist} (qui signifie « aider »).
\item Ne dites pas \eng{*I lost my time} : dites \eng{I \textbf{wasted} my time}.
\item Attention à l'ordre des mots : \eng{*I have well managed my time} est incorrect ; dites \eng{I managed my time \textbf{well}}.
\item \eng{Actually} signifie « en fait », pas « actuellement » ; pour « actuellement », dites \eng{currently}.
\item \eng{An event} prend l'article \eng{an} (voyelle) ; et l'on dit \eng{to arrive \textbf{at} the market}, jamais \eng{*to arrive to the market}.
\end{itemize}
\end{attention}

\begin{savaistu}[Le bénévolat dans le monde anglophone]
Dans les pays anglophones, le \eng{volunteering} (bénévolat) fait partie de la vie scolaire : au Ghana, au Nigeria ou aux États-Unis, les lycéens consacrent régulièrement des heures à leur communauté (\eng{community service}), et ces activités sont valorisées dans les dossiers d'admission à l'université. Au Cameroun anglophone, les \eng{work days} dans les écoles et les quartiers relèvent de la même tradition de solidarité.
\end{savaistu}

\courssub{Proposition de production et corrigé commenté}

\begin{exoresolu}[Modèle de composition commenté]
\textbf{Énoncé.} Rédiger une composition de 180 à 220 mots intitulée \emph{“Taking part in our Community Work Day and managing my time”}, en trois parties, avec au moins six connecteurs et cinq mots du vocabulaire de la leçon.
\tcblower
\textbf{Production modèle (environ 210 mots) :}

\emph{Taking part in our Community Work Day and managing my time}

Last Saturday, our English Club organised a Community Work Day in Yaoundé, and I was proud to take part in it. It was a wonderful experience of teamwork and commitment.

First of all, we met at Mokolo market at 7 a.m. sharp, because punctuality was our first rule. We cleaned the stalls and swept the ground together with the traders. Then, in the afternoon, we went to the municipal library, where we arranged books and read stories to the children. Moreover, I volunteered to lead the reading session, although I was a little nervous at first. Managing my time was not easy: I had to follow a strict timetable and I could not waste a single minute. However, thanks to good planning, everything went well. As a result, I understood that when you schedule your activities, you can serve your community and still have time for your studies.

In conclusion, this day taught me that taking part in community work is a duty and a pleasure. I am going to volunteer again next month, and I will manage my time even better.

\textbf{Commentaire des choix de langue :}
\begin{itemize}
\item \textbf{Structure} : introduction (présentation de l'événement), développement en deux paragraphes (actions puis gestion du temps), conclusion avec ouverture sur l'avenir.
\item \textbf{Connecteurs utilisés (8)} : \eng{first of all}, \eng{because}, \eng{then}, \eng{moreover}, \eng{although}, \eng{however}, \eng{as a result}, \eng{in conclusion}.
\item \textbf{Vocabulaire de la leçon (8 mots)} : \eng{to take part in}, \eng{teamwork}, \eng{commitment}, \eng{punctuality}, \eng{timetable}, \eng{to waste}, \eng{to volunteer}, \eng{to schedule}.
\item \textbf{Temps verbaux} : past simple pour le récit (\eng{met}, \eng{cleaned}, \eng{went}, \eng{taught}), present simple pour la vérité générale (\eng{is a duty and a pleasure}), futur pour le projet (\eng{am going to volunteer}, \eng{will manage}).
\item \textbf{Pièges évités} : \eng{took part in} (et non \emph{assisted to}), \eng{waste a minute} (et non \emph{lose my time}).
\end{itemize}
\end{exoresolu}

\begin{exercice}[Entraînement au Bac — sujet supplémentaire]
Rédigez, en 180 à 220 mots, une composition intitulée \emph{“How I take part in school activities and manage my time during the examination period”}. Contraintes : trois parties, au moins six connecteurs différents, cinq mots du vocabulaire de la leçon, temps verbaux cohérents. Temps conseillé : 45 minutes.
\end{exercice}

\begin{corrige}[Entraînement au Bac]
\textbf{Pistes de corrigé.} Introduction : présenter la période des examens et annoncer le sujet (\eng{During the examination period, I take part in many school activities, but I must manage my time carefully.}). Développement : (a) les activités (sport, English Club, devoirs) au present simple ; (b) la gestion du temps (\eng{timetable}, \eng{priorities}, \eng{I never waste time}) ; (c) une anecdote au past simple. Conclusion : bilan et projet au futur (\eng{Next term, I am going to...}). Connecteurs attendus : \eng{first of all}, \eng{moreover}, \eng{however}, \eng{because}, \eng{as a result}, \eng{in conclusion}. Vérifier : pas de \eng{*assist to}, pas de \eng{*lose time}, adverbe bien placé (\eng{I manage my time well}).
\end{corrige}

\courssub{Bilan de l'activité}

Cette activité d'intégration a permis de mobiliser l'ensemble des savoir-faire de la leçon dans une tâche authentique et motivante, directement inspirée de la vie scolaire camerounaise. L'élève a d'abord \textbf{compris} un document authentique (l'annonce du club), puis \textbf{planifié} sa production à l'aide d'un plan en trois parties, \textbf{rédigé} une composition complète en temps limité comme à l'examen, et enfin \textbf{évalué} et \textbf{corrigé} une production grâce à une grille critériée identique à celle du baccalauréat. L'évaluation croisée par binômes développe en outre l'esprit critique et la coopération : formuler \emph{two strengths and one piece of advice} apprend à lire un texte avec un regard de correcteur. La relecture finale, étape trop souvent négligée, est ici institutionnalisée : c'est elle qui transforme un brouillon en copie de qualité. Pour le Bac, retenez la formule gagnante : \cle{plan} + \cle{connecteurs} + \cle{vocabulaire précis} + \cle{relecture} = une composition réussie.

\newpage

\coursec{Méthodes et exercices résolus}

\coursec{Lesson 24 — Writing: Writes compositions on taking part and managing time in community/workplace/school activities}

\courssub{Savoir-faire 1 : Analyser le sujet et planifier sa composition}

\begin{methode}[Analyser le sujet et construire un plan en 5 minutes]
Avant d'écrire la moindre phrase, consacre 5 à 8 minutes à la préparation. C'est ce qui distingue une copie de 15/20 d'une copie de 8/20.
\begin{enumerate}
\item \textbf{Read the topic twice.} Souligne les mots-clés du sujet (\cle{discuss}, \cle{narrate}, \cle{describe}, \cle{give your opinion}) : ils déterminent le \textbf{type de composition} (narrative, descriptive, argumentative).
\item \textbf{Identify the constraints.} Qui écrit ? À qui ? Sur quel thème précis ? Ici : \emph{taking part} (participation) et \emph{managing time} (gestion du temps) dans des activités communautaires, scolaires ou professionnelles. Ne traite jamais un seul des deux axes si le sujet en demande deux.
\item \textbf{Brainstorm.} Note en vrac, en anglais, 6 à 8 idées : exemples vécus (a clean-up campaign in your quarter, a school football tournament, a community farm in Buea), problèmes de temps, solutions.
\item \textbf{Select and order.} Garde les 3 meilleures idées (une par paragraphe de développement) et classe-les du plus simple au plus fort.
\item \textbf{Draft a mini-plan} au brouillon : Introduction (hook + thesis) / Body 1, 2, 3 / Conclusion (summary + personal opinion). Chaque paragraphe = une idée + un exemple + un connecteur.
\end{enumerate}
\textbf{Réflexe Bac :} une composition de Terminale fait environ 250 à 300 mots. Un plan clair te fait gagner du temps au lieu d'en perdre.
\end{methode}

\begin{exoresolu}[Analyse de sujet type Bac]
\textbf{Énoncé.} Read the following topic and answer the questions.\par
\eng{“Students in your school complain that they have no time to take part in community activities. Write a composition in which you explain the importance of participating in such activities and suggest how students can manage their time better.”}\par
1. What type of composition is required (narrative, descriptive or argumentative)?\par
2. How many parts does the topic contain? Name them.\par
3. Propose a three-paragraph plan for the body of the composition.
\tcblower
\textbf{Solution détaillée.}\par
1. Le sujet demande d'\emph{expliquer} (explain) et de \emph{suggérer} (suggest) : il s'agit d'une composition \textbf{argumentative} (avec une dimension explicative), pas d'un récit. Le verbe \eng{explain} impose de donner des raisons ; \eng{suggest} impose de proposer des solutions.\par
2. Le sujet contient \textbf{deux parties} obligatoires : (a) \emph{the importance of participating in community activities} ; (b) \emph{how students can manage their time better}. Oublier l'une des deux = hors sujet partiel, très sanctionné.\par
3. Plan possible :\par
\textbf{Paragraph 1} — Why community activities matter: they build solidarity and teach responsibility (example: a clean-up campaign in the neighbourhood).\par
\textbf{Paragraph 2} — Benefits for the student himself: new skills, teamwork, a stronger CV (example: volunteering at a health centre in Yaoundé).\par
\textbf{Paragraph 3} — Time management solutions: making a weekly timetable, setting priorities, reducing time on the phone, choosing weekend activities.\par
\textbf{Conclusion :} chaque idée du plan correspond à un paragraphe ; l'introduction et la conclusion viendront encadrer ces trois paragraphes.
\end{exoresolu}

\courssub{Savoir-faire 2 : Structurer introduction, développement, conclusion}

\begin{definition}[Thesis statement]
La \cle{thesis statement} (phrase de thèse) est la phrase unique, placée à la fin de l'introduction, qui annonce le sujet précis et la position de l'auteur ou le plan de l'essai. Exemple : \eng{“This essay will show why such participation is valuable and how students can find time for it.”}
\end{definition}

\begin{propriete}[Structure type d'une composition argumentative]
\begin{center}
\begin{tabularx}{\linewidth}{|X|X|X|}
\hline
\rowcolor{popblueL} \textbf{Partie} & \textbf{Contenu} & \textbf{Longueur indicative} \\
\hline
Introduction & Hook (amorce) + Contextualisation + Thesis statement & 3-4 phrases \\
\hline
Development (Body) & 2 à 3 paragraphes : Topic sentence + Explication + Exemple + Transition & 4-5 phrases chacun \\
\hline
Conclusion & Résumé des points principaux (reformulé) + Opinion personnelle / Ouverture & 2-3 phrases \\
\hline
\end{tabularx}
\end{center}
\end{propriete}

\begin{methode}[Construire les trois parties de la composition]
\begin{enumerate}
\item \textbf{Introduction (3-4 phrases).} Commence par un \cle{hook} (une question, un constat général) : \eng{“In many Cameroonian towns, young people take part in community work every weekend.”} Puis annonce le sujet et ta \cle{thesis} (ta position ou ton plan) : \eng{“This essay will show why such participation is valuable and how students can find time for it.”}
\item \textbf{Development (2-3 paragraphes).} Un paragraphe = une \cle{topic sentence} (phrase d'idée) + une explication + un exemple concret + une phrase de transition. Jamais deux idées principales dans le même paragraphe.
\item \textbf{Conclusion (2-3 phrases).} Résume sans répéter mot pour mot, puis donne ton avis personnel ou une ouverture : \eng{“To conclude, community activities are not a waste of time but an investment, provided that students learn to organise their week.”}
\item \textbf{Paragraphing visible.} Saute une ligne ou fais un alinéa entre chaque partie : le correcteur doit \emph{voir} la structure.
\end{enumerate}
\end{methode}

\begin{exoresolu}[Remettre une composition dans l'ordre]
\textbf{Énoncé.} The following sentences are the jumbled parts of a composition. Number them from 1 to 5 in the correct order.\par
\eng{a) For instance, last month, I helped to clean the market square in my quarter on Saturday morning, after finishing my homework on Friday evening.}\par
\eng{b) In conclusion, taking part in community life is possible for every student who plans his week carefully.}\par
\eng{c) Many students believe that community activities take too much time; however, with good organisation, everyone can participate.}\par
\eng{d) First of all, participating in community work teaches us responsibility and solidarity.}\par
\eng{e) Moreover, managing our time well — for example by preparing a timetable — allows us to combine studies and community life.}
\tcblower
\textbf{Solution détaillée.}\par
Ordre correct : \textbf{c (1) — d (2) — a (3) — e (4) — b (5)}.\par
Justification : \eng{c} est une phrase d'introduction (constat général + thèse avec \eng{however}). \eng{d} ouvre le développement avec le connecteur de premier point \eng{First of all}. \eng{a} illustre cette première idée (\eng{For instance} = exemple concret, donc il suit l'idée qu'il illustre). \eng{e} ajoute le deuxième axe du sujet (la gestion du temps) grâce au connecteur d'addition \eng{Moreover}. \eng{b} conclut, signalé par \eng{In conclusion}.\par
\textbf{Leçon de méthode :} les connecteurs sont des panneaux indicateurs ; repère-les pour ordonner, et utilise-les pour écrire.
\end{exoresolu}

\courssub{Savoir-faire 3 : Utiliser les connecteurs logiques et les bons temps verbaux}

\begin{methode}[Choisir le bon connecteur et le bon temps]
\textbf{Connecteurs à mobiliser (apprends-en 2 par catégorie) :}
\begin{itemize}
\item Ordonner : \cle{first of all}, \cle{then}, \cle{finally}
\item Ajouter : \cle{moreover}, \cle{in addition}, \cle{furthermore}
\item Opposer : \cle{however}, \cle{whereas}, \cle{on the contrary}
\item Donner une cause : \cle{because}, \cle{since}, \cle{as}
\item Donner une conséquence : \cle{therefore}, \cle{as a result}, \cle{consequently}
\item Illustrer : \cle{for instance}, \cle{for example}
\item Conclure : \cle{to conclude}, \cle{in conclusion}, \cle{all in all}
\end{itemize}
\textbf{Temps verbaux :} vérités générales et habitudes $\rightarrow$ \textbf{present simple} (\eng{Students take part in clean-up campaigns.}) ; expérience passée racontée $\rightarrow$ \textbf{past simple} (\eng{Last year, I joined the school health club.}) ; projets et intentions $\rightarrow$ \textbf{going to / will} (\eng{We are going to organise a car wash next month.}).
\end{methode}

\begin{exoresolu}[Compléter avec connecteurs et verbes conjugués]
\textbf{Énoncé.} Complete the paragraph with the correct connector from the box and the verb in brackets at the right tense.\par
Box : \eng{however — moreover — as a result — for instance}\par
\eng{Many workers in Douala say they are too busy for community activities. \ldots (1), a well-organised person can always find one or two hours a week. Last year, my uncle \ldots (2) (join) the neighbourhood development association. He \ldots (3) (attend) meetings every Sunday after church. \ldots (4), he also coaches the youth football team. \ldots (5), he has become one of the most respected people in the quarter.}
\tcblower
\textbf{Solution détaillée.}\par
(1) \eng{However} : opposition entre « trop occupés » et « on peut trouver du temps ». Majuscule en début de phrase.\par
(2) \eng{joined} : \eng{last year} impose le \textbf{past simple} ; verbe régulier, \emph{join + ed}.\par
(3) \eng{attends} : \eng{every Sunday} exprime une habitude présente $\rightarrow$ \textbf{present simple}, 3\textsuperscript{e} personne du singulier, d'où le \textbf{-s}. Attention : ne pas écrire \emph{attend} sans -s.\par
(4) \eng{Moreover} : ajout d'une activité supplémentaire (coaching).\par
(5) \eng{As a result} : conséquence logique de son engagement (il est respecté).\par
\textbf{Paragraphe corrigé :} \eng{“Many workers in Douala say they are too busy for community activities. However, a well-organised person can always find one or two hours a week. Last year, my uncle joined the neighbourhood development association. He attends meetings every Sunday after church. Moreover, he also coaches the youth football team. As a result, he has become one of the most respected people in the quarter.”}
\end{exoresolu}

\courssub{Savoir-faire 4 : Employer le vocabulaire de la participation et de la gestion du temps}

\begin{methode}[Activer le lexique thématique]
Une copie de Bac doit contenir du \textbf{vocabulaire précis} du thème. Prépare deux colonnes mentales :
\begin{itemize}
\item \textbf{Participation :} \cle{to take part in}, \cle{to participate in}, \cle{to volunteer}, \cle{to join a club / an association}, \cle{community work}, \cle{a clean-up campaign}, \cle{teamwork}, \cle{a sense of responsibility}, \cle{to contribute to}.
\item \textbf{Gestion du temps :} \cle{to manage one's time}, \cle{time management}, \cle{a timetable / a schedule}, \cle{to set priorities}, \cle{to plan ahead}, \cle{to waste time}, \cle{to meet a deadline}, \cle{to balance studies and activities}.
\end{itemize}
\textbf{Réflexe :} varie les verbes (\eng{take part in / participate in / join / get involved in}) pour éviter les répétitions, et relie chaque mot de vocabulaire à un exemple concret du contexte camerounais (a \emph{njangi} group, a village development committee, a school environmental club).
\end{methode}

\begin{savaistu}[Le « Community Work » au Cameroun]
Au Cameroun, la participation communautaire est une réalité quotidienne : les \emph{development associations} (comités de développement villageois ou de quartier) organisent des \emph{clean-up campaigns} (opérations « ville propre ») le samedi matin, la construction de ponts, l'entretien des routes rurales. Les \emph{njangi} (tontines solidaires) financent des projets collectifs. À l'école, les \emph{environmental clubs} et les \emph{health clubs} mènent des actions de sensibilisation. Le 11 février (Fête de la Jeunesse) voit une mobilisation massive des élèves pour le défilé et les travaux d'intérêt général. Intégrer ces réalités dans ta composition montre ton ancrage culturel et impressionne le correcteur.
\end{savaistu}

\begin{exoresolu}[Vocabulaire : traduction et réemploi]
\textbf{Énoncé.} a) Translate into English : « Les élèves devraient participer aux activités communautaires et apprendre à gérer leur temps. »\par
b) Complete : \eng{If you want to succeed, you must \ldots your priorities and stop \ldots time on your phone.} (set / waste, aux formes correctes)\par
c) Find in the lesson a synonym of \eng{to take part in}.
\tcblower
\textbf{Solution détaillée.}\par
a) \eng{“Students should take part in community activities and learn to manage their time.”} Justifications : \eng{should} pour le conseil (devraient) ; \eng{take part in} + nom (participer \emph{à}) ; \eng{community activities} — attention, \emph{communautaires} ne se traduit pas par \emph{communitarian} ; \eng{to manage their time} — \emph{gérer} se dit \emph{manage}, pas \emph{control}.\par
b) \eng{“If you want to succeed, you must set your priorities and stop wasting time on your phone.”} Après \eng{must}, base verbale : \eng{set}. Après \eng{stop} au sens de « cesser de », le \textbf{gérondif} : \eng{wasting} (et non \emph{to waste}, qui changerait le sens : « s'arrêter pour gaspiller »).\par
c) \eng{to participate in} ou \eng{to get involved in}.\par
\textbf{Conclusion :} le vocabulaire juste + la construction juste (préposition, gérondif) = les deux tiers de la note d'expression écrite.
\end{exoresolu}

\courssub{Savoir-faire 5 : Relire et corriger sa production}

\begin{methode}[La relecture en 4 passes (5 minutes)]
\begin{enumerate}
\item \textbf{Passe « contenu » :} ai-je traité \emph{toutes} les parties du sujet ? Chaque paragraphe a-t-il une idée claire et un exemple ?
\item \textbf{Passe « grammaire » :} accords sujet-verbe (\eng{he takes}, pas \emph{he take}), temps cohérents, -ing après les prépositions (\eng{good at organising}).
\item \textbf{Passe « orthographe et ponctuation » :} majuscules en début de phrase et au pronom \eng{I}, pas de virgule entre sujet et verbe, apostrophes (\eng{students'} au pluriel possessif).
\item \textbf{Passe « présentation » :} alinéas visibles, écriture lisible, nombre de mots respecté (environ 250-300).
\end{enumerate}
\textbf{Réflexe :} corrige d'abord les erreurs que tu fais \emph{toujours} (tiens une liste personnelle de tes fautes fréquentes).
\end{methode}

\begin{exoresolu}[Corriger un extrait de copie d'élève]
\textbf{Énoncé.} The following extract contains 5 mistakes. Find and correct them.\par
\eng{“Last saturday, I have participated in a clean-up campaign in my quarter. We was very tired but happy. The community activities learns us solidarity. I think that manage time is very important for students.”}
\tcblower
\textbf{Solution détaillée.}\par
1. \emph{saturday} $\rightarrow$ \eng{Saturday} : les jours prennent la \textbf{majuscule} en anglais.\par
2. \emph{have participated} $\rightarrow$ \eng{participated} : \eng{last Saturday} est un repère temporel \textbf{terminé} $\rightarrow$ past simple, jamais de present perfect avec un repère passé défini (calque du passé composé français).\par
3. \emph{We was} $\rightarrow$ \eng{We were} : le past simple de \emph{be} est \emph{were} avec \emph{we}.\par
4. \emph{learns us} $\rightarrow$ \eng{teach us} : faux ami — « apprendre quelque chose à quelqu'un » = \eng{to teach somebody something} ; \emph{learn} = apprendre (recevoir le savoir). Accord : \eng{activities} (pluriel) $\rightarrow$ \eng{teach}, sans -s ! Correction complète : \eng{Community activities teach us solidarity.}\par
5. \emph{manage time is} $\rightarrow$ \eng{managing time is} : un verbe sujet se met au \textbf{gérondif} en anglais (calque de l'infinitif sujet français « gérer son temps est… »).\par
\textbf{Phrase corrigée :} \eng{“Last Saturday, I participated in a clean-up campaign in my quarter. We were very tired but happy. Community activities teach us solidarity. I think that managing time is very important for students.”}
\end{exoresolu}

\courssub{Savoir-faire 6 : Rédiger une composition complète en conditions d'examen}

\begin{methode}[De la consigne à la copie en 45 minutes]
\begin{enumerate}
\item \textbf{5 min :} analyse du sujet + mini-plan au brouillon (voir savoir-faire 1).
\item \textbf{5 min :} introduction au brouillon (hook + thesis).
\item \textbf{25 min :} rédaction au propre, un paragraphe à la fois, en suivant le plan ; vise 3 paragraphes de développement de 4-5 phrases chacun.
\item \textbf{5 min :} conclusion (résumé + opinion personnelle).
\item \textbf{5 min :} relecture en 4 passes (voir savoir-faire 5).
\end{enumerate}
\textbf{Grille d'évaluation type Bac (à connaître) :} respect du sujet et du type de texte (4 pts) ; organisation et connecteurs (4 pts) ; correction grammaticale (4 pts) ; richesse du vocabulaire (4 pts) ; orthographe et présentation (4 pts).
\end{methode}

\begin{exoresolu}[Composition type Bac — modèle commenté]
\textbf{Énoncé.} \eng{Your school is launching a “Community Service Week”. Write a composition (about 250 words) explaining why students should take part in it and how they can manage their time to combine studies and community activities.}
\tcblower
\textbf{Corrigé rédigé (modèle) :}\par
\eng{“Nowadays, many students spend all their free time on their phones and forget the world around them. The Community Service Week organised by our school is an excellent opportunity to change this habit. This essay will explain why every student should participate and how good time management makes it possible.}\par
\eng{First of all, community activities teach values that books cannot give. When students clean a market square or visit an orphanage, they learn solidarity, humility and responsibility. For instance, last year, my classmates and I repainted the walls of a primary school in our quarter, and the children's smiles taught us more than any lesson.}\par
\eng{Moreover, taking part in such activities is useful for the students themselves. Teamwork develops communication skills, and voluntary work is always appreciated on a CV or in a job interview.}\par
\eng{However, many students claim they have no time. The solution is simple: they must manage their week better. Preparing a timetable, setting priorities and reducing time spent on social media can easily free two or three hours. As a result, studies and community life can be balanced without stress.}\par
\eng{To conclude, community service is not a waste of time but an investment in ourselves and in our society. If every student plans his week carefully, the Community Service Week will be a great success.”}\par
\textbf{Commentaire :} introduction avec hook + thesis annoncée ; 3 paragraphes de développement (valeurs / bénéfices personnels / gestion du temps) ouverts par \eng{First of all / Moreover / However} ; exemple concret camerounais ; conclusion avec opinion personnelle. Environ 230-250 mots : conforme à l'attendu.
\end{exoresolu}

\begin{attention}[Les 5 erreurs qui coûtent des points au Bac]
\begin{enumerate}
\item \textbf{Present perfect avec un repère passé :} \emph{I have participated last year.} Faux ! Avec \eng{last year}, \eng{yesterday}, \eng{ago} $\rightarrow$ past simple : \eng{I participated last year.} Calque du passé composé français.
\item \textbf{Faux amis :} « apprendre à qqn » = \eng{teach} (pas \emph{learn}) ; « assister à » = \eng{attend} (pas \emph{assist}, qui signifie « aider ») ; « une activité » = \eng{activity} (pluriel \eng{activities}, pas \emph{activitie}) ; « gérer » = \eng{manage} (pas \emph{control}).
\item \textbf{Infinitif sujet :} \emph{Manage time is important.} Faux ! En anglais, le verbe sujet est au gérondif : \eng{Managing time is important.}
\item \textbf{Oubli du -s de la 3\textsuperscript{e} personne et des majuscules :} \emph{he take part}, \emph{on saturday}. Correct : \eng{He takes part}, \eng{on Saturday}. Jours, mois et le pronom \eng{I} prennent toujours la majuscule.
\item \textbf{Prépositions et constructions :} \emph{participate to} $\rightarrow$ \eng{participate in} ; \emph{good to organise} $\rightarrow$ \eng{good at organising} ; \emph{stop to waste time} $\rightarrow$ \eng{stop wasting time}. Toute préposition est suivie du gérondif.
\end{enumerate}
\end{attention}

\newpage

\coursec{Exercices}

\courssub{Je vérifie mes connaissances}

\begin{exercice}[QCM : vocabulaire et collocations]
Pour chaque phrase, choisis la bonne réponse (A, B ou C).

\begin{enumerate}
\item Students should \cle{save} time by planning their week. \\
A. make time \quad B. save time \quad C. do time
\item Don't \ldots\ your time watching television all day! \\
A. win \quad B. spend \quad C. waste
\item We must \ldots\ an effort to participate in community work. \\
A. do \quad B. make \quad C. take
\item The headmaster asked me to \ldots\ the meeting on Friday. \\
A. assist \quad B. attend \quad C. participate
\item Our school \ldots\ a clean-up campaign every month. \\
A. organizes \quad B. does \quad C. makes
\item I am very busy, but I always find time \ldots\ my homework. \\
A. for doing \quad B. to do \quad C. doing
\end{enumerate}
\end{exercice}

\begin{exercice}[Vrai ou faux — justifie ta réponse]
Dis si les affirmations suivantes sont vraies (V) ou fausses (F), puis justifie en une phrase en français.

\begin{enumerate}
\item A composition usually has three parts: an introduction, a body and a conclusion.
\item The connector \eng{however} is used to add a similar idea.
\item \eng{First of all}, \eng{then} and \eng{finally} are used to order ideas.
\item In a composition, you can start writing immediately without a plan.
\item \eng{In conclusion} introduces the final paragraph of a composition.
\item The verb \eng{to participate} is followed by the preposition \eng{to}.
\end{enumerate}
\end{exercice}

\begin{exercice}[Vocabulaire : participation et gestion du temps]
Complète chaque phrase avec le mot anglais qui convient, choisi dans la liste suivante :

\begin{center}
\eng{deadline — schedule — volunteer — punctual — take part in — priority — community — manage}
\end{center}

\begin{enumerate}
\item Many students \ldots\ the cleaning of their school every Friday.
\item A \ldots\ is a person who works for free to help others.
\item You must respect the \ldots\ : the composition must be handed in on Monday.
\item My \ldots\ is very busy: lessons from 7 a.m. to 5 p.m., then football training.
\item To succeed, you must learn to \ldots\ your time well.
\item Being \ldots\ means arriving on time, never late.
\item Finishing my homework is my first \ldots\ .
\item Our \ldots\ organised a meeting to build a new well in the village.
\end{enumerate}
\end{exercice}

\begin{exercice}[Conjugaison : les bons temps]
Mets le verbe entre parenthèses à la forme et au temps qui conviennent.

\begin{enumerate}
\item Every Saturday, we \ldots\ (clean) the market square with our neighbours.
\item Last year, my school \ldots\ (organise) a big sports competition.
\item If students \ldots\ (manage) their time well, they will pass their exams.
\item I \ldots\ (never / waste) my time since I started using a timetable.
\item Next term, we \ldots\ (take part) in a tree-planting campaign in Buea.
\item While the teacher \ldots\ (explain) the lesson, some students were chatting.
\end{enumerate}
\end{exercice}

\courssub{Je m'entraîne}

\begin{exercice}[Traduction]
Traduis les phrases suivantes en anglais.

\begin{enumerate}
\item Les élèves doivent apprendre à gérer leur temps s'ils veulent réussir.
\item Tout d'abord, participer aux activités de la communauté développe le sens des responsabilités.
\item De plus, le bénévolat permet aux jeunes d'acquérir de nouvelles compétences.
\item Cependant, beaucoup d'élèves perdent leur temps sur les réseaux sociaux.
\item En conclusion, une bonne organisation est la clé du succès scolaire.
\end{enumerate}
\end{exercice}

\begin{exercice}[Texte à trous : les connecteurs]
Complète le texte suivant avec les connecteurs de la liste (chaque connecteur ne peut être utilisé qu'une seule fois) :

\begin{center}
\eng{first of all — however — moreover — as a result — in conclusion}
\end{center}

\begin{textebox}[Managing time at school]
Time is the most precious gift a student has. \ldots\ , every student should prepare a weekly timetable and follow it strictly. \ldots\ , it is important to avoid wasting time on useless activities such as playing games all day. \ldots\ , many students say they have no time to study; this is not true: they simply do not organise themselves. A student who plans his day well finishes his homework early and \ldots\ , he has free time to rest and to help his family. \ldots\ , managing time well is the secret of success at school and in life.
\end{textebox}
\end{exercice}

\begin{exercice}[Corrige les erreurs]
Les huit phrases suivantes contiennent chacune une erreur typique des francophones. Souligne l'erreur et réécris la phrase correctement.

\begin{enumerate}
\item I have seventeen years and I am in Lower Sixth.
\item Many students participate to the clean-up campaign.
\item He manage his time very well.
\item Actually, I am preparing my Bac examination.
\item I am agree with the idea of volunteering.
\item She is a very good student because she works hardly.
\item We must to respect the deadlines.
\item The teacher gave us many advices about time management.
\end{enumerate}
\end{exercice}

\begin{exercice}[Appariement : notes et topic sentences]
Voici des notes de brainstorming (colonne A) et des phrases de sujet (\emph{topic sentences}) possibles (colonne B). Associe chaque série de notes à la \emph{topic sentence} qui convient, puis indique si le paragraphe irait dans le développement ou la conclusion d'une composition sur \eng{Taking part in community activities}.

\begin{center}
\begin{tabularx}{\linewidth}{|X|X|}
\hline
\rowcolor{popblueL} \textbf{A — Notes} & \textbf{B — Topic sentences} \\
\hline
1. clean the streets, plant trees, build wells & a. To sum up, community work benefits both the individual and society. \\
\hline
2. learn new skills, make friends, become responsible & b. Community activities help young people to grow personally. \\
\hline
3. a better village, a cleaner town, a stronger nation & c. There are many ways in which students can serve their community. \\
\hline
\end{tabularx}
\end{center}
\end{exercice}

\courssub{Je me prépare au Bac}

\begin{exercice}[Compréhension et langue — type Bac]
Lis attentivement le texte suivant, puis réponds aux questions en anglais, par des phrases complètes.

\begin{textebox}[A lesson in time management]
When Ngo Bassa entered the Lower Sixth at Government High School, Yaoundé, she was always tired and her marks were falling. She attended lessons all day, helped her mother at the market in the evening, and still wanted to play handball with her friends. “I have no time to study,” she often complained.

One day, her English teacher, Mr Fombah, gave the class a simple exercise: “Write down everything you do from Monday to Sunday, hour by hour.” Ngo Bassa was surprised by the result. She discovered that she spent three hours a day on her phone and two hours watching television. “The problem is not time,” Mr Fombah told her, “the problem is how you use it.”

Ngo Bassa decided to change. She prepared a weekly timetable: homework first, then housework, then sport, and only thirty minutes of telephone in the evening. She also joined the school's environmental club, which cleans the classrooms every Friday afternoon. Little by little, her marks improved, and she felt less tired. At the end of the year, she was among the best students of her class.

“Managing my time changed my life,” she now tells the younger students. “If I can do it, everybody can.”
\end{textebox}

\textbf{Comprehension questions}
\begin{enumerate}
\item Why was Ngo Bassa always tired at the beginning of the school year? (2 marks)
\item What exercise did Mr Fombah give to the class? (2 marks)
\item What did Ngo Bassa discover about herself? (2 marks)
\item Mention two changes she made in her daily life. (2 marks)
\item What lesson can we learn from Ngo Bassa's story? (2 marks)
\end{enumerate}

\textbf{Language work}
\begin{enumerate}
\item Find in the text the opposite of: a) \eng{worse} b) \eng{refused}. (2 marks)
\item Complete with the right tense: Ngo Bassa \ldots\ (improve) her marks since she started using a timetable. (1 mark)
\item Rewrite in the reported speech: “I have no time to study,” she often complained. (2 marks)
\item Give a suitable title to the text and justify your choice in one sentence. (1 mark)
\end{enumerate}
\end{exercice}

\begin{exercice}[Traduction — type Bac]
Traduis en anglais le passage suivant.

\begin{textebox}[Texte à traduire]
Dans mon village, les jeunes participent activement aux travaux communautaires. Chaque premier samedi du mois, nous nettoyons le marché et réparons les routes. Grâce à ces activités, nous apprenons à gérer notre temps et à travailler en équipe. De plus, les anciens nous transmettent leurs connaissances. Je pense que chaque élève devrait consacrer quelques heures par semaine au service de sa communauté, car c'est ainsi que l'on construit un avenir meilleur.
\end{textebox}

\emph{Barème indicatif : respect du sens (4 points), exactitude grammaticale (4 points), richesse du vocabulaire (2 points).}
\end{exercice}

\begin{exercice}[Composition — type Bac]
\textbf{Topic:} \emph{The importance of managing time for a student.}

Write a composition of about \textbf{200 words} on the topic above. You have \textbf{45 minutes}.

Your composition must:
\begin{itemize}
\item have a clear \cle{introduction} (with a hook and a plan), two developed \cle{body paragraphs} and a \cle{conclusion};
\item use at least \textbf{five} logical connectors (e.g. \eng{first of all}, \eng{moreover}, \eng{however}, \eng{as a result}, \eng{in conclusion});
\item include examples from school or community life in Cameroon;
\item use appropriate tenses and vocabulary related to time management and participation.
\end{itemize}

\emph{Grille d'évaluation : contenu et pertinence (4 pts), organisation et connecteurs (4 pts), correction grammaticale (4 pts), vocabulaire (4 pts), présentation générale (4 pts). Total : 20 points.}
\end{exercice}

\newpage

\coursec{Corrigés des exercices}

\begin{corrige}[Exercice 1 — QCM : vocabulaire et collocations]
\begin{enumerate}
\item \textbf{B. save time.} \eng{Students should save time by planning their week.} La collocation correcte est \eng{to save time} (gagner du temps). \eng{Make time} signifie « trouver le temps » et \eng{do time} signifie « faire de la prison ».
\item \textbf{C. waste.} \eng{Don't waste your time watching television all day!} \eng{To waste time} = perdre son temps. \eng{Spend time} est neutre ou positif ; \eng{win time} n'existe pas en anglais (faux ami du français « gagner du temps » au sens de perdre du temps).
\item \textbf{B. make.} \eng{We must make an effort to participate in community work.} La collocation est \eng{to make an effort} ; on ne dit jamais \eng{do an effort} ni \eng{take an effort}.
\item \textbf{B. attend.} \eng{The headmaster asked me to attend the meeting on Friday.} \eng{To attend a meeting} = assister à une réunion. Attention : \eng{assist} signifie « aider » (faux ami), et \eng{participate} exige la préposition \eng{in}.
\item \textbf{A. organizes.} \eng{Our school organizes a clean-up campaign every month.} On \eng{organizes} une campagne ; \eng{do} et \eng{make} ne conviennent pas avec \eng{campaign}.
\item \textbf{B. to do.} \eng{I always find time to do my homework.} Après \eng{find time}, on utilise l'infinitif avec \eng{to} : \eng{find time to do something}.
\end{enumerate}
\end{corrige}

\begin{corrige}[Exercice 2 — Vrai ou faux]
\begin{enumerate}
\item \textbf{V.} Une composition anglaise classique comprend trois parties : \eng{introduction}, \eng{body} (développement) et \eng{conclusion}.
\item \textbf{F.} \eng{However} exprime l'opposition ou le contraste (« cependant ») ; pour ajouter une idée semblable, on utilise \eng{moreover}, \eng{furthermore} ou \eng{in addition}.
\item \textbf{V.} \eng{First of all}, \eng{then} et \eng{finally} sont des connecteurs d'ordre chronologique qui organisent la progression des idées.
\item \textbf{F.} Il faut toujours planifier avant de rédiger : brainstorming, plan en trois parties, choix des connecteurs ; écrire sans plan produit un texte désorganisé.
\item \textbf{V.} \eng{In conclusion} (ou \eng{to conclude}, \eng{to sum up}) introduit le paragraphe final qui résume l'argumentation.
\item \textbf{F.} \eng{To participate} est suivi de la préposition \eng{in} : \eng{to participate in an activity}. C'est une erreur typique du francophone (« participer à »).
\end{enumerate}
\end{corrige}

\begin{corrige}[Exercice 3 — Vocabulaire : participation et gestion du temps]
\begin{enumerate}
\item \eng{Many students \textbf{take part in} the cleaning of their school every Friday.} (participer à)
\item \eng{A \textbf{volunteer} is a person who works for free to help others.} (un bénévole)
\item \eng{You must respect the \textbf{deadline}: the composition must be handed in on Monday.} (la date limite)
\item \eng{My \textbf{schedule} is very busy: lessons from 7 a.m. to 5 p.m., then football training.} (l'emploi du temps)
\item \eng{To succeed, you must learn to \textbf{manage} your time well.} (gérer)
\item \eng{Being \textbf{punctual} means arriving on time, never late.} (ponctuel)
\item \eng{Finishing my homework is my first \textbf{priority}.} (priorité)
\item \eng{Our \textbf{community} organised a meeting to build a new well in the village.} (la communauté)
\end{enumerate}
\emph{Point de vigilance :} \eng{schedule} se prononce « ché-dioul » en anglais britannique et « ské-dioul » en anglais américain.
\end{corrige}

\begin{corrige}[Exercice 4 — Conjugaison : les bons temps]
\begin{enumerate}
\item \eng{Every Saturday, we \textbf{clean} the market square with our neighbours.} Présent simple : habitude signalée par \eng{every Saturday}.
\item \eng{Last year, my school \textbf{organised} a big sports competition.} Prétérit simple : action passée datée (\eng{last year}). Verbe régulier : \emph{organise, organised, organised}.
\item \eng{If students \textbf{manage} their time well, they will pass their exams.} Après \eng{if} (condition), présent simple dans la subordonnée, futur dans la principale. Jamais de futur après \eng{if}.
\item \eng{I \textbf{have never wasted} my time since I started using a timetable.} Present perfect : action commencée dans le passé qui continue (\eng{since}). Attention à la place de \eng{never} : entre l'auxiliaire et le participe passé.
\item \eng{Next term, we \textbf{will take part} (ou \textbf{are going to take part}) in a tree-planting campaign in Buea.} Futur : projet annoncé par \eng{next term}.
\item \eng{While the teacher \textbf{was explaining} the lesson, some students were chatting.} Prétérit continu après \eng{while} : action longue en arrière-plan.
\end{enumerate}
\end{corrige}

\begin{corrige}[Exercice 5 — Traduction]
\begin{enumerate}
\item \eng{Students must learn to manage their time if they want to succeed.} — « gérer » = \eng{to manage}, jamais \eng{to control} ni \eng{to drive}.
\item \eng{First of all, taking part in community activities develops the sense of responsibility.} — Sujet en \emph{-ing} (\eng{taking part}) ; « sens des responsabilités » = \eng{sense of responsibility}.
\item \eng{Moreover, volunteering allows young people to acquire new skills.} — \eng{to allow somebody to do something} ; « compétences » = \eng{skills}.
\item \eng{However, many students waste their time on social networks.} — « perdre son temps » = \eng{to waste one's time}, pas \eng{to lose one's time}.
\item \eng{In conclusion, good organisation is the key to success at school.} — On dit \eng{the key to success} (préposition \eng{to}).
\end{enumerate}
\end{corrige}

\begin{corrige}[Exercice 6 — Texte à trous : les connecteurs]
\begin{enumerate}
\item \eng{\textbf{First of all}, every student should prepare a weekly timetable and follow it strictly.} (première idée du développement)
\item \eng{\textbf{Moreover}, it is important to avoid wasting time on useless activities such as playing games all day.} (ajout d'une idée)
\item \eng{\textbf{However}, many students say they have no time to study; this is not true: they simply do not organise themselves.} (opposition)
\item \eng{A student who plans his day well finishes his homework early and \textbf{as a result}, he has free time to rest and to help his family.} (conséquence)
\item \eng{\textbf{In conclusion}, managing time well is the secret of success at school and in life.} (conclusion)
\end{enumerate}
\emph{Justification :} chaque connecteur a une fonction logique précise — ordre (\eng{first of all}), addition (\eng{moreover}), contraste (\eng{however}), résultat (\eng{as a result}), synthèse (\eng{in conclusion}). Les intervertir détruit la cohérence du texte.
\end{corrige}

\begin{corrige}[Exercice 7 — Corrige les erreurs]
\begin{enumerate}
\item \eng{I \textbf{am seventeen years old} and I am in Lower Sixth.} — En anglais, on \emph{est} son âge (\eng{to be}), on ne l'\emph{a} pas ; et \eng{years old} est obligatoire.
\item \eng{Many students participate \textbf{in} the clean-up campaign.} — \eng{To participate in}, pas \eng{to} (calque du français « participer à »).
\item \eng{He \textbf{manages} his time very well.} — Au présent simple, le verbe prend un \emph{-s} à la troisième personne du singulier.
\item \eng{\textbf{At present} (ou \textbf{Currently}), I am preparing my Bac examination.} — \eng{Actually} signifie « en fait » ; c'est un faux ami de « actuellement ».
\item \eng{I \textbf{agree} with the idea of volunteering.} — \eng{Agree} est un verbe : on dit \eng{I agree}, jamais \eng{I am agree}.
\item \eng{She is a very good student because she works \textbf{hard}.} — \eng{Hard} est à la fois adjectif et adverbe ; \eng{hardly} signifie « à peine ».
\item \eng{We must \textbf{respect} the deadlines.} — Après un modal (\eng{must}, \eng{can}, \eng{should}), l'infinitif est sans \eng{to}.
\item \eng{The teacher gave us \textbf{a lot of advice} about time management.} — \eng{Advice} est indénombrable : pas de pluriel en \emph{-s} ; on dit \eng{some advice}, \eng{a piece of advice}.
\end{enumerate}
\end{corrige}

\begin{corrige}[Exercice 8 — Appariement : notes et topic sentences]
\begin{itemize}
\item \textbf{1 — c.} \eng{There are many ways in which students can serve their community.} Les notes (nettoyer les rues, planter des arbres, construire des puits) énumèrent des \emph{actions concrètes} : c'est un paragraphe du \textbf{développement} (premier paragraphe).
\item \textbf{2 — b.} \eng{Community activities help young people to grow personally.} Les notes (acquérir des compétences, se faire des amis, devenir responsable) décrivent les \emph{bénéfices personnels} : paragraphe du \textbf{développement} (deuxième paragraphe).
\item \textbf{3 — a.} \eng{To sum up, community work benefits both the individual and society.} Les notes (un meilleur village, une ville plus propre, une nation plus forte) généralisent et concluent ; le connecteur \eng{to sum up} confirme qu'il s'agit de la \textbf{conclusion}.
\end{itemize}
\emph{Méthode :} une \eng{topic sentence} résume l'idée générale du paragraphe ; les notes sont les exemples qui la développent. La présence d'un connecteur de synthèse (\eng{to sum up}) signale toujours une conclusion.
\end{corrige}

\begin{corrige}[Exercice 9 — Compréhension et langue — type Bac]
\textbf{Comprehension questions — réponses attendues (barème : 2 points par réponse complète et exacte)}
\begin{enumerate}
\item \eng{Ngo Bassa was always tired because she attended lessons all day, helped her mother at the market in the evening and still wanted to play handball, so she had no time to rest and study.}
\item \eng{Mr Fombah asked the class to write down everything they did from Monday to Sunday, hour by hour.}
\item \eng{She discovered that she wasted a lot of time: she spent three hours a day on her phone and two hours watching television.}
\item \eng{She prepared a weekly timetable with clear priorities (homework first, then housework, then sport) and she reduced her telephone time to thirty minutes in the evening. She also joined the school's environmental club.} (deux changements suffisent)
\item \eng{We learn that the problem is not the lack of time but the way we use it: good time management leads to success.}
\end{enumerate}

\textbf{Language work}
\begin{enumerate}
\item a) \eng{worse} $\rightarrow$ \eng{better} ; b) \eng{refused} $\rightarrow$ \eng{accepted} (ou \eng{decided to change} $\rightarrow$ l'antonyme direct attendu est \eng{accepted}). (1 point par réponse)
\item \eng{Ngo Bassa \textbf{has improved} her marks since she started using a timetable.} Present perfect obligatoire avec \eng{since}. (1 point)
\item \eng{She often complained that she had no time to study.} Discours indirect : \eng{have} $\rightarrow$ \eng{had} (recul du temps), suppression des guillemets, introduction par \eng{that}. (2 points)
\item Titre possible : \eng{“Time Well Managed Is Success Earned”} — justification : \eng{This title is suitable because the text shows how Ngo Bassa succeeded by learning to manage her time.} (1 point)
\end{enumerate}

\emph{Points de vigilance :} répondre par des phrases complètes (sujet + verbe) ; reprendre le temps de la question ; au discours indirect, ne jamais garder le présent après un verbe introducteur au prétérit.
\end{corrige}

\begin{corrige}[Exercice 10 — Traduction — type Bac]
\textbf{Proposition de traduction modèle :}

\eng{In my village, young people take an active part in community work. Every first Saturday of the month, we clean the market and repair the roads. Thanks to these activities, we learn to manage our time and to work as a team. Moreover, the elders pass on their knowledge to us. I think that every student should devote a few hours a week to serving his or her community, because that is how we build a better future.}

\textbf{Barème indicatif (10 points) :}
\begin{itemize}
\item Respect du sens (4 pts) : toutes les idées du texte français sont rendues, sans omission ni ajout.
\item Exactitude grammaticale (4 pts) : concordance des temps (présents d'habitude), prépositions correctes (\eng{take part in}, \eng{thanks to}, \eng{pass on \ldots{} to}), articles.
\item Richesse du vocabulaire (2 pts) : \eng{community work}, \eng{as a team}, \eng{devote \ldots{} to}, \eng{build a better future}.
\end{itemize}

\emph{Points de vigilance :} « les anciens » = \eng{the elders}, pas \eng{the ancients} ; « consacrer du temps à » = \eng{to devote time to} + nom ou \emph{-ing} ; « c'est ainsi que » = \eng{that is how}, pas \eng{it is like that that}.
\end{corrige}

\begin{corrige}[Exercice 11 — Composition — type Bac]
\textbf{Proposition de rédaction modèle (environ 200 mots) :}

\eng{Time is the most precious resource a student possesses, yet many young people waste it without realising it. Managing time well is therefore essential for success at school and in life.}

\eng{First of all, a student who manages his time well works better. By preparing a weekly timetable and doing his homework before watching television, he studies regularly and avoids the stress of last-minute work. As a result, his marks improve and he feels more confident.}

\eng{Moreover, good time management leaves room for other important activities. A well-organised student can take part in community life: in my school in Yaoundé, for example, the environmental club cleans the classrooms every Friday afternoon, and its members still find time to study and to play football. However, students who spend hours on social networks complain that they have no time; the truth is that they simply do not organise themselves.}

\eng{In conclusion, managing time is the key to academic success and to a balanced life. Every student should learn this lesson early, because time lost can never be recovered.}

\textbf{Grille d'évaluation (20 points) :}
\begin{itemize}
\item Contenu et pertinence (4 pts) : le sujet est traité, les arguments sont développés avec des exemples camerounais.
\item Organisation et connecteurs (4 pts) : introduction avec accroche, deux paragraphes de développement, conclusion ; au moins cinq connecteurs (\eng{first of all}, \eng{moreover}, \eng{as a result}, \eng{however}, \eng{in conclusion}).
\item Correction grammaticale (4 pts) : temps corrects (présent de vérité générale), accords, prépositions.
\item Vocabulaire (4 pts) : lexique de la gestion du temps et de la participation (\eng{timetable}, \eng{waste time}, \eng{take part in}, \eng{well-organised}).
\item Présentation générale (4 pts) : lisibilité, paragraphes nets, respect du nombre de mots.
\end{itemize}

\emph{Points de vigilance :} ne pas recopier le sujet comme introduction ; annoncer un plan dans l'introduction ; un exemple concret par paragraphe ; relire les terminaisons en \emph{-s} à la troisième personne et les faux amis (\eng{actually}, \eng{assist}, \eng{advices}).
\end{corrige}

\newpage

\coursec{Fiche bilan}

\begin{popfigure}
\begin{tikzpicture}[mindmap, grow cyclic, every node/.style={concept}, text width=2.6cm, align=center,
  root concept/.append style={concept color=popblue, font=\bfseries, text=white, text width=3cm},
  level 1/.append style={level distance=3.6cm, sibling angle=51.4},
  level 1 concept/.append style={font=\small\bfseries}]
\node[root concept]{Writing: taking part \& managing time}
  child[concept color=poporange]{ node {Analyse du sujet \& planification} }
  child[concept color=popgreen]{ node {Structure: intro, body, conclusion} }
  child[concept color=poppurple]{ node {Connecteurs logiques} }
  child[concept color=poppink]{ node {Temps verbaux} }
  child[concept color=popgold]{ node {Vocabulaire: participation \& time} }
  child[concept color=popblue!60]{ node {Modèle \& erreurs à éviter} }
  child[concept color=popgreen!60]{ node {Grille d'évaluation Bac} };
\end{tikzpicture}
\legende{Carte mentale de la leçon : rédiger une composition sur la participation et la gestion du temps.}
\end{popfigure}

\begin{bilan}
\textbf{1. Lexique clé (à connaître par cœur)}
\begin{itemize}
\item \eng{to take part in} = participer à ; \eng{to get involved in} = s'impliquer dans ; \eng{to volunteer} = faire du bénévolat.
\item \eng{community service} = service communautaire ; \eng{teamwork} = travail d'équipe ; \eng{a duty / a responsibility} = un devoir / une responsabilité.
\item \eng{to manage one's time} = gérer son temps ; \eng{a timetable / a schedule} = un emploi du temps ; \eng{a deadline} = une date limite.
\item \eng{to set priorities} = fixer des priorités ; \eng{to waste time} = perdre son temps ; \eng{time-consuming} = chronophage.
\item \eng{punctual} = ponctuel ; \eng{committed} = engagé ; \eng{reliable} = fiable ; \eng{efficient} = efficace.
\end{itemize}

\textbf{2. Grammaire à maîtriser}
\begin{center}
\begin{tabularx}{\linewidth}{|X|X|}
\hline
\rowcolor{popblueL} Règle & Exemple \\
\hline
Prétérit simple : actions passées accomplies & \eng{Last year, I joined the school choir.} \\
\hline
Present perfect : expérience, bilan au présent & \eng{I have learnt to manage my time better.} \\
\hline
Présent simple : habitudes, vérités générales & \eng{Students who plan their week succeed.} \\
\hline
Gérondif après \eng{to be used to}, \eng{to spend time}, prépositions & \eng{I spend two hours cleaning the school.} \\
\hline
Modaux : \eng{should, must, can} pour conseils et obligations & \eng{We should respect deadlines.} \\
\hline
\end{tabularx}
\end{center}

\textbf{3. Connecteurs logiques indispensables}
\begin{itemize}
\item Addition : \eng{first of all, moreover, furthermore, in addition}.
\item Opposition : \eng{however, nevertheless, whereas, on the other hand}.
\item Cause et conséquence : \eng{because, since, therefore, as a result, consequently}.
\item Ordre et temps : \eng{first, then, next, afterwards, finally}.
\item Conclusion : \eng{to conclude, in short, all in all, to sum up}.
\end{itemize}

\textbf{4. Méthode express en 5 étapes}
\begin{enumerate}
\item \cle{Analyser} le sujet : souligner les mots-clés, dégager la problématique.
\item \cle{Brainstormer} : noter idées et vocabulaire en 5 minutes.
\item \cle{Planifier} : introduction (accroche + annonce du plan), 2 ou 3 paragraphes de développement (une idée + un exemple chacun), conclusion (bilan + ouverture).
\item \cle{Rédiger} en liant les idées avec des connecteurs et en contrôlant les temps.
\item \cle{Relire} : accords, temps, orthographe, ponctuation, nombre de mots.
\end{enumerate}
\end{bilan}

\begin{aretenir}[Checklist avant le Bac]
Avant de rendre ta copie, vérifie chaque point :
\begin{itemize}
\item $\square$ J'ai bien compris le sujet et je réponds exactement à la question posée.
\item $\square$ Ma composition a une introduction, un développement (2 ou 3 paragraphes) et une conclusion.
\item $\square$ Mon introduction contient une accroche et annonce clairement le plan.
\item $\square$ Chaque paragraphe développe une seule idée principale avec un exemple (école, communauté, travail).
\item $\square$ J'ai utilisé au moins cinq connecteurs logiques variés et correctement placés.
\item $\square$ J'ai employé le vocabulaire du thème : \eng{to take part in, to manage one's time, deadline, priorities, volunteer}.
\item $\square$ Mes temps verbaux sont cohérents (prétérit pour le passé, présent pour les habitudes).
\item $\square$ J'ai évité les faux amis : \eng{actually} = en fait, \eng{to assist} n'est pas « assister à » (dire \eng{to attend}).
\item $\square$ J'ai vérifié la troisième personne du singulier (\eng{he manages}, \eng{she takes part}).
\item $\square$ J'ai relu l'orthographe et la ponctuation (majuscule en début de phrase, point final).
\item $\square$ J'ai respecté le nombre de mots demandé et présenté une copie lisible.
\item $\square$ Ma conclusion apporte un bilan personnel et une ouverture, sans répéter l'introduction.
\end{itemize}
\end{aretenir}

\findecours

\end{document}
